\subsection{将函数解释为缓增广义函数}

定义 5. --- 设 \(\nu\) 是 \(\mathbf{R}^{n}\) 上的测度。如果存在整数 \(r \in \mathbf{N}\) ，使得连续映射 \(x \mapsto (1 + \|x\|)^{-r}\) 关于 \(\nu\) 在 \(\mathbf{R}^{n}\) 上可积，就称它是缓增的。

换言之，测度 \(\nu\) 是缓增的，如果存在 \(r \in \mathbf{N}\) ，使得由 \(x \mapsto \| x\|^{-r}\) 所定义的函数关于 \(\nu\) 在 \(\mathbf{R}^n\) 中单位球的余集上可积。特别地， \(\mathbf{R}^n\) 上的每个有界测度都是缓增的。更一般地，若 \(f\) 是多项式增长的 \(\mu\)-可测函数，且 \(\nu\) 是缓增的，则测度 \(f \cdot \nu\) 是缓增的。

\(\mathcal{M}^{\mathrm{t}}(\mathbf{R}^{n})\) （\(\mathbf{R}^{n}\) 上缓增测度所成的集）是 \(\mathcal{M}(\mathbf{R}^{n};\mathbf{C})\) （\(\mathbf{R}^{n}\) 上复测度所成的空间）的一个向量子空间。

命题 15. --- 设 \(\nu\) 是 \(\mathbf{R}^{n}\) 上的缓增测度。 \(\nu\) 在 \(\mathcal{S}(\mathbf{R}^{n})\) 上的限制是一个缓增广义函数。它为零当且仅当测度 \(\nu\) 为零。

由于 \(\nu\) 是缓增的，存在正整数 \(k\) ，使得映射 \(x \mapsto \| x\|^{-k}\) 关于 \(\nu\) 在 \(\mathbf{R}^n\) 中单位球的余集上可积。对每个施瓦茨函数 \(\varphi \in \mathcal{S}(\mathbf{R}^n)\) ，我们有

\[
| \langle \nu , \varphi \rangle | \leqslant \Bigl (\int_ {\| x \| \leqslant 1} d \nu \Bigr) \widetilde {q} _ {0, 0} (\varphi) + \Bigl (\int_ {\| x \| > 1} \| x \| ^ {- k} d \nu \Bigr) \widetilde {q} _ {0, k} (\varphi),
\]

因此 \(\nu\) 在 \(\mathcal{S}(\mathbf{R}^n)\) 上的限制是一个缓增广义函数。

最后的断言由 IV, p. 206 的命题 6 得出，因为 \(\mathcal{D}(\mathbf{R}^n)\) 包含在 \(\mathcal{S}(\mathbf{R}^n)\) 中。

我们将把空间 \(\mathcal{M}^{\mathrm{t}}(\mathbf{R}^{n})\) 等同于 \(\mathcal{S}'(\mathbf{R}^{n})\) 的一个子空间。

命题 16. --- 设 \(p \in [1, +\infty]\) ， \(f \in \mathcal{L}^{p}(\mathbf{R}^{n})\) 。那么测度 \(f \cdot \mu\) （以 \(f\) 为密度，关于勒贝格测度）是缓增的。映射 \(f \mapsto f \cdot \mu\) （从 \(\mathrm{L}^{p}(\mathbf{R}^{n})\) 到 \(\mathcal{S}'(\mathbf{R}^{n})\) 中的）是连续的。

设 \(q\) 是 \(p\) 的共轭指数，并设 \(r \geqslant 0\) 满足 \(rq > n\) 。对每个 \(x \in \mathbf{R}^n\) ，令 \(g(x) = (1 + \|x\|)^{-r}\) 。函数 \(g\) 属于 \(\mathcal{L}^q(\mathbf{R}^n)\) （由 IV, p. 199 的命题 3）。由赫尔德不等式，我们有

\[
\int_ {\mathbf {R} ^ {n}} ^ {*} (1 + \| x \|) ^ {- r} | f (x) | d \mu (x) \leqslant \mathrm{N} _ {q} (g) \mathrm{N} _ {p} (f) <   + \infty ,
\]

因此测度 \(f \cdot \mu\) 是缓增的。

设 \(f \in \mathrm{L}^p(\mathbf{R}^n)\) ， \(\varphi \in \mathcal{S}(\mathbf{R}^n)\) 。赫尔德不等式蕴含

\[
| \langle f \cdot \mu , \varphi \rangle | = \left| \int_ {\mathbf {R} ^ {n}} f (x) \varphi (x) d \mu (x) \right| \leqslant \| f \| _ {p} \| \varphi \| _ {q}
\]

而映射 \(f \mapsto f \cdot \mu\) 的连续性则由 IV, p. 213 的命题 13 得出。

对每个 \(p \in [1, +\infty]\) ，我们将把 \(\mathrm{L}^p(\mathbf{R}^n)\) 等同于 \(\mathcal{S}'(\mathbf{R}^n)\) 的一个子空间（借助线性映射 \(f \mapsto f \cdot \mu\) ）。

命题 17. --- 空间 \(\mathcal{D}(\mathbf{R}^{n})\) 与 \(\mathcal{S}(\mathbf{R}^{n})\) 在 \(\mathcal{S}'(\mathbf{R}^{n})\) 中稠密。

只须证明 \(\mathcal{D}(\mathbf{R}^{n})\) 在 \(\mathcal{S}'(\mathbf{R}^{n})\) 中稠密。设 \(\lambda\) 是 \(\mathcal{S}'(\mathbf{R}^{n})\) 上的连续线性形式，它在 \(\mathcal{D}(\mathbf{R}^{n})\) 上为零。由于空间 \(\mathcal{S}(\mathbf{R}^{n})\) 是自反的，存在函数 \(\varphi \in \mathcal{S}(\mathbf{R}^{n})\) ，使得对所有 \(\lambda(f) = \langle f, \varphi \rangle\) 有 \(f \in \mathcal{S}'(\mathbf{R}^{n})\) 。于是我们有

\[
0 = \lambda (\psi) = \langle \psi , \varphi \rangle = \int_ {\mathbf {R} ^ {n}} \psi \varphi d \mu
\]

对所有 \(\psi\in\mathcal{D}(\mathbf{R}^{n})\) 成立。因此 \(\varphi\cdot\mu\) 上的测度 \(\mathbf{R}^{n}\) 为零（IV, p. 206 的命题 6），从而 \(\varphi=0\) 。命题于是由哈恩–巴拿赫定理（EVT, II, p. 46，推论 1）得出。

\subsection{缓增广义函数的傅里叶变换}

由于每个施瓦茨函数 \(\varphi\) 在 \(\mathbf{R}^n\) 上可积（IV, p. 213 的命题 13），它有傅里叶变换 \(\mathcal{F}(\varphi)\) （相应地，余傅里叶变换 \(\overline{\mathcal{F}}(\varphi)\) ），后者等同于 \(\mathbf{R}^n\) 上由下式定义的连续有界函数

\[
y \mapsto \int_ {\mathbf {R} ^ {n}} \varphi (x) \exp (- 2 i \pi x \cdot y) d \mu (x)
\]

（相应地，函数 \(y \mapsto \int_{\mathbf{R}^n} \varphi(x) \exp(2i\pi x \cdot y) d\mu(x)\) ）。

引理 11. --- 设 \(\varphi\in\mathcal{S}(\mathbf{R}^{n})\) 。函数 \(\mathcal{F}(\varphi)\) 在 \(\mathbf{R}^{n}\) 上是无穷次可微的，且我们有

\[
\mathcal {F} (\mathrm{X} ^ {\alpha} \varphi) = (- 2 i \pi) ^ {- | \alpha |} \partial^ {\alpha} (\mathcal {F} (\varphi)),
\]

\[
\mathcal {F} (\partial^ {\alpha} \varphi) = (2 i \pi) ^ {| \alpha |} X ^ {\alpha} \mathcal {F} (\varphi)
\]

对所有 \(\alpha\) （属于 \(\mathbf{N}^n\) ）成立。

不妨设 \(n \geqslant 1\) 。设 \(\varphi \in \mathcal{S}(\mathbf{R}^{n})\) 。由 \((x, y) \mapsto \varphi(x) \exp(2i\pi x \cdot y)\) （从 \(\mathbf{R}^{n} \times \mathbf{R}^{n}\) 到 \(\mathbf{C}\) 中的函数）对每个整数 k 满足 IV, p. 198 的推论 1 的假设。因此 \(\varphi\) 的傅里叶变换是无穷次可微的，且满足

\[
\partial^ {\alpha} (\mathcal {F} \varphi) (y) = (- 2 i \pi) ^ {| \alpha |} \int_ {\mathbf {R} ^ {n}} x ^ {\alpha} \varphi (x) \exp (- 2 i \pi x \cdot y) d \mu (x)
\]

对所有 \(y \in \mathbf{R}^{n}\) 成立，由此即得第一个公式。

我们来证明第二个公式。对 \(|\alpha|\) 作归纳，只须处理 \(|\alpha| = 1\) 的情形，而易化为 \(\alpha = (1,0,\ldots,0)\) 的情形。把每个 \(x \in \mathbf{R}^n\) 写成 \(x = (x_1,x')\) 的形式，其中 \(x' \in \mathbf{R}^{n-1}\) ，并用 \(\mu_1\) （相应地， \(\mu'\) ）表示 \(\mathbf{R}^{n-1}\) （相应地， \(\mathbf{R}\) ）上的勒贝格测度。由勒贝格–富比尼定理（INT, V, p. 96, § 8, no. 4，定理 1），对每个 \(y = (y_1,y') \in \mathbf{R} \times \mathbf{R}^{n-1}\) ，得到

\[
\begin{array}{l} \mathcal {F} (\partial_ {1} \varphi) (y) = \int_ {\mathbf {R} ^ {n - 1}} \exp (- 2 i \pi   x ^ {\prime} \cdot y ^ {\prime}) \\ \qquad \times \Big (\int_ {\mathbf {R}} (\partial_ {1} \varphi) (x _ {1}, x ^ {\prime}) \exp (- 2 i \pi   x _ {1} y _ {1}) d \mu_ {1} (x _ {1}) \Big) d \mu^ {\prime} (x ^ {\prime}). \end{array}
\]

对每个紧区间 \([a,b]\) （R 中的）以及每个 \(x'\in R^{n-1}\) ，我们有

\[
\begin{array}{r l} & {\int_ {a} ^ {b} (\partial_ {1} \varphi) (x _ {1}, x ^ {\prime}) \exp (- 2 i \pi x _ {1} y _ {1}) d \mu_ {1} (x _ {1}) =} \\ & {\qquad \left[ \varphi (x _ {1}, x ^ {\prime}) \exp (- 2 i \pi x _ {1} y _ {1}) \right] _ {a} ^ {b}} \\ & {\qquad + 2 i \pi y _ {1} \int_ {a} ^ {b} \varphi (x _ {1}, x ^ {\prime}) \exp (- 2 i \pi x _ {1} y _ {1}) d \mu_ {1} (x _ {1})} \end{array}
\]

由分部积分（FVR, II, p. 10，公式 (10)）。当 \(a\) 趋于 \(-\infty\) 且 \(b\) 趋于 \(+\infty\) 时，右端第一项趋于 0，因为 \(\varphi \in \mathcal{S}(\mathbf{R}^n)\) 。第二项由勒贝格定理（INT, IV, p. 137, § 3, no. 7，定理 6）收敛于

\[
2 i \pi y _ {1} \int_ {\mathbf {R}} \varphi (x _ {1}, x ^ {\prime}) \exp (- 2 i \pi x _ {1} y _ {1}) d \mu_ {1} (x _ {1}),
\]

因为映射 \(x_{1} \mapsto \varphi(x_{1}, x')\) 在 R 上可积。由于 \(x_{1} \mapsto \partial_{1}\varphi(x_{1}, x')\) 也在 R 上可积，由此得出

\[
\begin{array}{r l r} & & {\int_ {\mathbf {R}} \partial_ {1} \varphi (x _ {1}, x ^ {\prime}) \exp (- 2 i \pi x _ {1} y _ {1}) d \mu_ {1} (x _ {1})} \\ & & {= 2 i \pi y _ {1} \int_ {\mathbf {R}} \varphi (x _ {1}, x ^ {\prime}) \exp (- 2 i \pi x _ {1} y _ {1}) d \mu_ {1} (x _ {1})} \end{array}
\]

最后，再次应用勒贝格–富比尼定理，我们得出

\[
\mathcal {F} (\partial_ {1} \varphi) (y) = 2 i \pi y _ {1} \int_ {\mathbf {R} ^ {n}} \varphi (x) \exp (- 2 i \pi x \cdot y) d \mu (x),
\]

这正是所要证明的。

命题 18. --- 傅里叶变换在 \(\mathcal{S}(\mathbf{R}^{n})\) 上的限制是拓扑向量空间的一个自同构，其逆是余傅里叶变换的限制。

设 \(\varphi\in\mathcal{S}(\mathbf{R}^{n})\) 。由前一引理， \(\varphi\) 的傅里叶变换属于 \(\mathcal{C}^{\infty}(\mathbf{R}^{n})\) 。此外，对 \(\alpha\in\mathbf{N}^{n}\) 与 \(\beta\in\mathbf{N}^{n}\) ，我们有

\[
\begin{array}{r} \mathrm{X} ^ {\beta} \partial^ {\alpha} (\mathcal {F} (\varphi)) = (- 2 i \pi) ^ {| \alpha |} \mathrm{X} ^ {\beta} \mathcal {F} (\mathrm{X} ^ {\alpha} \varphi) \\ = (- 1) ^ {| \alpha |} (2 i \pi) ^ {| \alpha | - | \beta |} \mathcal {F} (\partial^ {\beta} (\mathrm{X} ^ {\alpha} \varphi)). \end{array}
\]

特别地，函数 \(\mathrm{X}^{\beta}\partial^{\alpha}(\mathcal{F}(\varphi))\) 有界。由于 \(\alpha\) 与 \(\beta\) 在 \(\mathbf{N}^{n}\) 中任意，这表明 \(\mathcal{F}(\varphi)\) 属于 \(\mathcal{S}(\mathbf{R}^{n})\) 。此外，这一计算蕴含

\[
q _ {\alpha , \beta} (\mathcal {F} (\varphi)) \leqslant (2 \pi) ^ {| \alpha | - | \beta |} \| \partial^ {\beta} (X ^ {\alpha} \varphi) \| _ {1}
\]

对 \((\alpha, \beta) \in \mathbf{N}^n \times \mathbf{N}^n\) 与 \(\varphi \in \mathcal{S}(\mathbf{R}^n)\) 成立。

由于空间 \(\mathcal{S}(\mathbf{R}^{n})\) 在 \(\mathrm{L}^{1}(\mathbf{R}^{n})\) 中的包含映射是连续的（IV, p. 213 的命题 13），映射 \(q_{\alpha,\beta} \circ \mathcal{F}\) （从 \(\mathcal{S}(\mathbf{R}^{n})\) 到 \(\mathbf{R}\) 中的）是连续的。由此得出，傅里叶变换是从空间 \(\mathcal{S}(\mathbf{R}^{n})\) 到其自身的连续映射（参见 EVT, II, p. 7，命题 5, c)）。类似地可以验证，余傅里叶变换是从空间 \(\mathcal{S}(\mathbf{R}^{n})\) 到其自身的连续映射。由傅里叶反演公式（II, p. 222 的定理 3），傅里叶变换与余傅里叶变换是互逆的同构。

定义 6. --- \(\mathcal{S}'(\mathbf{R}^n)\) 上的傅里叶变换（相应地，余傅里叶变换）称为 \(\mathcal{S}(\mathbf{R}^n)\) 上的傅里叶变换（相应地，余傅里叶变换）的转置。

我们仍用 \(\mathcal{F}\) （相应地， \(\overline{\mathcal{F}}\) ）表示 \(\mathcal{S}'(\mathbf{R}^n)\) 上的傅里叶变换（相应地，余傅里叶变换）。 \(\mathcal{S}'(\mathbf{R}^n)\) 上的傅里叶变换于是是一个拓扑向量空间自同构，其逆是余傅里叶变换。对每个 \(f \in \mathcal{S}'(\mathbf{R}^n)\) ，缓增广义函数 \(\mathcal{F}(f)\) （相应地， \(\overline{\mathcal{F}}(f) )\) ）由下述公式定义

\[
\langle \mathcal {F} (f), \varphi \rangle = \langle f, \mathcal {F} (\varphi) \rangle \quad (\text { resp. } \langle \overline {{\mathcal {F}}} (f), \varphi \rangle = \langle f, \overline {{\mathcal {F}}} (\varphi) \rangle)
\]

对所有 \(\varphi\in\mathcal{S}(\mathbf{R}^{n})\) 成立。

命题 19. --- 设 f 是与有界测度 \(\nu \in \mathcal{M}^{1}(\mathbf{R}^{n})\) （相应地，与 \(g \in \mathrm{L}^{2}(\mathbf{R}^{n})\) ）相伴的缓增广义函数。 f 在 \(\mathcal{S}'(\mathbf{R}^{n})\) 中的傅里叶变换是与测度 \(\nu\) 的傅里叶变换（相应地，与 g 的傅里叶变换）相伴的缓增广义函数。

设 \(\nu\) 是 \(\mathbf{R}^{n}\) 上的有界测度， f 是与 \(\nu\) 相伴的缓增广义函数。傅里叶变换 \(\mathcal{F}(\nu)\) 是 \(\mathbf{R}^{n}\) 上的连续有界函数（II, p. 207 的命题 3）。与这一函数相伴的缓增广义函数满足

\[
\langle \mathcal {F} (\nu), \varphi \rangle = \int_ {\mathbf {R} ^ {n}} \mathcal {F} (\nu) \varphi d \mu = \int_ {\mathbf {R} ^ {n}} \mathcal {F} (\varphi) d \nu = \langle f, \mathcal {F} (\varphi) \rangle = \langle \mathcal {F} (f), \varphi \rangle
\]

对所有 \(\varphi\in\mathcal{S}(\mathbf{R}^{n})\) 成立，其中第二个等式就是 II, p. 221 的命题 13，由于测度 \(\varphi\cdot\mu\) 有界，它在此适用。因此与 \(\mathcal{F}(\nu)\) 相伴的缓增广义函数是 \(\mathcal{F}(f)\) 。

当 \(f\) 是与 \(g \in \mathrm{L}^2(\mathbf{R}^n)\) 相伴的缓增广义函数时，我们用完全类似的方法，借助 II, p. 221 的公式 (29) 处理。

对余傅里叶变换也有类似的命题。

注. --- 关于测度的傅里叶变换的初等公式对缓增广义函数的傅里叶变换仍然有效。例如，对 \(f \in \mathcal{S}'(\mathbf{R}^{n})\) 与 \(\alpha \in \mathbf{N}^{n}\) ，我们有

\[
\begin{array}{l} \mathcal {F} (\partial^ {\alpha} f) = (2 i \pi) ^ {| \alpha |} X ^ {\alpha} \mathcal {F} (f) \\ \mathcal {F} (X ^ {\alpha} f) = (- 2 i \pi) ^ {- | \alpha |} \partial^ {\alpha} (\mathcal {F} (f)) \end{array}
\]

这由引理 11 得出，

\subsection{向量空间上的广义函数与缓增广义函数}

设 \(u\) 是从 \(\mathbf{R}^{n}\) 到 \(\mathbf{R}^{n}\) 中的双射线性映射。映射 \(\varphi \mapsto \varphi \circ u\) 是 \(\mathcal{S}(\mathbf{R}^{n})\) （相应地， \(\mathcal{D}(\mathbf{R}^{n})\) ）的自同构；它的转置是 \(\mathcal{S}'(\mathbf{R}^{n})\) （相应地， \(\mathcal{D}'(\mathbf{R}^{n})\) ）的自同构。

设 E 是 n 维实向量空间。设 \(v: \mathbf{R}^{n} \to E\) 是向量空间同构。我们用 \(\mathcal{S}(E)\) （相应地， \(\mathcal{D}(E)\) ）表示映射 \(\varphi: E \to \mathbf{C}\) 中满足 \(\varphi \circ v \in \mathcal{S}(\mathbf{R}^{n})\) （相应地，满足 \(\varphi \circ v \in \mathcal{D}(\mathbf{R}^{n})\) ）者所成的集。由前面的注，这些空间不依赖于 v 的选取；它们同构于 \(\mathcal{S}(\mathbf{R}^{n})\) （相应地， \(\mathcal{D}(\mathbf{R}^{n})\) ）。我们用 \(\mathcal{S}'(E)\) （相应地， \(\mathcal{D}'(E)\) ）表示 \(\mathcal{S}(E)\) （相应地， \(\mathcal{D}(E)\) ）的对偶，赋有有界收敛拓扑。它是同构于 \(\mathcal{S}'(\mathbf{R}^{n})\) （相应地，同构于 \(\mathcal{D}'(\mathbf{R}^{n})\) ）的拓扑向量空间。

设 E 与 F 是 n 维实向量空间，它们关于双线性形式 \(b: E \times F \to \mathbf{R}\) 处于对偶。局部紧交换群 E 关于映射 

\[
(x, y) \mapsto \exp (2 i \pi b (x, y))
\]

（从 E × F 到 U 中；参见 II, p. 235 的推论 1）与 F 处于对偶。我们给 E 与 F 赋以关于这一映射互为对偶的哈尔测度。

空间 \(\mathcal{S}(\mathrm{E})\) 包含在 \(\mathrm{L}^1 (\mathrm{E})\) 中； E 的傅里叶变换通过过渡到子空间与对偶，诱导出从 \(\mathcal{S}(\mathrm{E})\) 到 \(\mathcal{S}(\mathrm{F})\) 中的拓扑向量空间同构，其转置是从 \(\mathcal{S}'(\mathrm{F})\) 到 \(\mathcal{S}'(\mathrm{E})\) 中的拓扑向量空间同构。

\subsection{索伯列夫空间}

设 U 是 \(\mathbf{R}^{n} \) 的开子集。设 p 是实数 \(\geqslant 1\) ， k 是自然数。我们用 \(W^{k,p}(U)\) 表示满足如下条件的广义函数 \(f \in \mathcal{D}'(U)\) 所成的空间：对每个 \(\alpha \in N^{n}\) （满足 \(|\alpha| \leqslant k\) ），广义函数 \(\partial^{\alpha}f\) 与 \(\mathrm{L}^{p}(U) \) 的一个元素相伴。

特别地，当 \(U = \mathbf{R}^{n}\) 时， \(W^{k,p}(U)\) 中的元素是缓增广义函数。

从 \(\mathrm{W}^{k,p}(\mathrm{U})\) 到 \(\mathbf{R}_{+}\) 中的、把 \(f\in \mathrm{W}^{k,p}(\mathrm{U})\)

\[
\| f \| _ {k, p} = \Bigl (\sum_ {| \alpha | \leqslant k} \| \partial^ {\alpha} f \| _ {p} ^ {p} \Bigr) ^ {1 / p}
\]

是 \(W^{k,p}(U)\) 上的一个范数。空间 \(W^{k,p}(U)\) 将始终赋有这一范数；这一赋范空间称为指标为 k 、指数为 p 的索伯列夫空间。

空间 \(\mathcal{D}(\mathrm{U})\) 包含在 \(\mathrm{W}^{k,p}(\mathrm{U})\) 中。我们用 \(\mathrm{W}_0^{k,p}(\mathrm{U})\) 表示 \(\mathcal{D}(\mathrm{U})\) 在 \(\mathrm{W}^{k,p}(\mathrm{U})\) 中的闭包。它是 \(\mathrm{W}^{k,p}(\mathrm{U})\) 的一个闭子空间。

有 \(\mathrm{W}_{0}^{k,p}(\mathbf{R}^{n})=\mathrm{W}^{k,p}(\mathbf{R}^{n})\) ，但空间 \(\mathrm{W}^{k,p}(\mathrm{U})\) 与 \(\mathrm{W}_{0}^{k,p}(\mathrm{U})\) 一般是不同的（参见 IV, p. 334 的习题 12 与 IV, p. 334 的习题 14）。

我们也记 \(\mathrm{H}^k (\mathrm{U}) = \mathrm{W}^{k,2}(\mathrm{U})\) ， \(\mathrm{H}_0^k (\mathrm{U}) = \mathrm{W}_0^{k,2}(\mathrm{U})\) 。

 \(\mathrm{H}^{k}(\mathrm{U})\) 的范数是准希尔伯特范数，它由 \(\mathrm{H}^{k}(\mathrm{U})\) 上的、由下式定义的正埃尔米特形式所导出。

\[
\left(f _ {1}, f _ {2}\right) \mapsto \sum_ {| \alpha | \leqslant k} \int_ {\mathrm{U}} \overline {{\partial^ {\alpha} f _ {1}}} \partial^ {\alpha} f _ {2} d \mu .
\]

希尔伯特空间 \(\mathrm{H}^k (\mathrm{U})\) 与 EVT, V, p. 6 例 (3) 中记作 \(\mathcal{H}^k\) 的空间重合。

按定义有 \(\mathrm{W}^{0,p}(\mathrm{U})=\mathrm{L}^{p}(\mathrm{U})\) ， \(\mathrm{H}^{0}(\mathrm{U})=\mathrm{L}^{2}(\mathrm{U})\) ；此外，由 IV, p. 202 的命题 4, b)， \(\mathrm{W}_{0}^{0,p}(\mathrm{U})=\mathrm{L}^{p}(\mathrm{U})\) 。

命题 20. — 索伯列夫空间 \(W^{k,p}(U)\) 与 \(W_{0}^{k,p}(U)\) 是可数型的巴拿赫空间。特别地，空间 \(H^{k}(U)\) 与 \(H_{0}^{k}(U)\) 是可数型的希尔伯特空间。

只须证明关于 \(\mathrm{W}^{k,p}(\mathrm{U})\) 的断言。

设 I 是满足 \(\alpha \in \mathbf{N}^n\) 的 \(|\alpha| \leqslant k\) 所成的集。线性映射 \(u\) （从 \(\mathrm{W}^{k,p}(\mathrm{U})\) 到 \(\mathrm{L}^p (\mathrm{U})^{\mathrm{I}}\) 中的）把 \(f\) 对应到族 \((\partial^{\alpha}f)_{\alpha \in \mathrm{I}}\) ，它是单射；按 \(\mathrm{W}^{k,p}(\mathrm{U})\) 上范数的定义，它是连续且严格的。为证 \(\mathrm{W}^{k,p}(\mathrm{U})\) 完备，只须证它在 \(u\) 下的像是闭的。设 \((f_n)_{n \in \mathbf{N}}\) 是 \(\mathrm{W}^{k,p}(\mathrm{U})\) 中的序列，使得 \((u(f_n))_{n \in \mathbf{N}}\) 收敛。设 \((g_{\alpha})_{\alpha \in \mathrm{I}} \in \mathrm{L}^p (\mathrm{U})^{\mathrm{I}}\) 是它的极限。对 \(\alpha \in \mathrm{I}\) ，序列 \((\partial^{\alpha}f_n)_{n \in \mathbf{N}}\) 在 \(\mathrm{L}^p (\mathrm{U})\) 中收敛于 \(g_{\alpha}\) 。当然，收敛也在 \(\mathcal{D}'(\mathrm{U})\) 中发生。令 \(f = g_0\) 。对每个 \(\varphi \in \mathcal{D}(\mathrm{U})\) ，得到

\[
\begin{array}{r l} & {\langle \partial^ {\alpha} f, \varphi \rangle = (- 1) ^ {| \alpha |} \langle f, \partial^ {\alpha} \varphi \rangle} \\ & {\qquad = \lim _ {n \to + \infty} (- 1) ^ {| \alpha |} \langle f _ {n}, \partial^ {\alpha} \varphi \rangle = \lim _ {n \to + \infty} \langle \partial^ {\alpha} f _ {n}, \varphi \rangle = \langle g _ {\alpha}, \varphi \rangle .} \end{array}
\]

这就证明了 \(g_{\alpha} = \partial^{\alpha} f\) 对所有 \(\alpha \in I\) 成立，从而 \((g_{\alpha})_{\alpha \in I} = u(f)\) 属于 u 的像。

空间 \(\mathrm{W}^{k,p}(\mathrm{U})\) 由 u 等同于空间 \(\mathrm{L}^{p}(\mathrm{U})^{\mathrm{I}}\) 的一个子空间；后者是可数型的（IV, p. 180 的命题 2 与 TG, IX, p. 19，cor., (ii))），因而 \(\mathrm{W}^{k,p}(\mathrm{U})\) 也是可数型的（loc. cit., (i)）。

命题 21. --- 设 N 是 \(\mathbf{R}^{n}\) 上的欧几里得范数。设 \(k\) 是整数 \(\geqslant 0\) 。索伯列夫空间 \(\mathrm{H}^{k}(\mathbf{R}^{n})\) 是这样的 \(f \in \mathcal{S}'(\mathbf{R}^{n})\) 所成的空间： \((1 + \mathrm{N}^{k})\mathcal{F}(f)\) 属于 \(\mathrm{L}^{2}(\mathbf{R}^{n})\) 。

我们对 \(k\) 作归纳。当 \(k = 0\) 时，结论由普朗谢雷尔定理（II, p. 215，定理 1）得出。设 \(k = 1\) 。对 \(f \in \mathcal{S}'(\mathbf{R}^n)\) ，我们有 \((1 + \mathrm{N})\mathcal{F}f \in \mathrm{L}^2(\mathbf{R}^n)\) 当且仅当 \(f \in \mathrm{L}^2(\mathbf{R}^n)\) 且 \(\mathrm{N}\mathcal{F}f \in \mathrm{L}^2(\mathbf{R}^n)\) 。此外， \(\mathrm{N}\mathcal{F} \in \mathrm{L}^2(\mathbf{R}^n)\) 当且仅当：对每个 \(\alpha \in \mathbf{N}^n\) （满足 \(|\alpha| = 1\) ），有 \(\mathrm{X}^\alpha\mathcal{F}f \in \mathrm{L}^2(\mathbf{R}^n)\) 。由于 \(\mathcal{F}(\partial^\alpha f) = 2i\pi\mathrm{X}^\alpha\mathcal{F}f\) ，这一条件意味着对每个 \(\mathcal{F}(\partial^\alpha f) \in \mathrm{L}^2(\mathbf{R}^n)\) （满足 \(\alpha\) ）有 \(|\alpha| = 1\) ，亦即对每个 \(\partial^\alpha f \in \mathrm{L}^2(\mathbf{R}^n)\) （满足 \(\alpha\) ）有 \(|\alpha| = 1\) 。由此得出断言对 \(k = 1\) 为真。

现在设 \(k \geqslant 2\) ，并且关于 \(\mathrm{H}^{\ell}(\mathbf{R}^{n})\) 的断言对每个正整数 \(\ell \leqslant k - 1\) 成立。设 \(f \in \mathcal{S}'(\mathbf{R}^{n})\) 。按定义， \(f \in \mathrm{H}^{k}(\mathbf{R}^{n})\) 当且仅当 \(f \in \mathrm{L}^{2}(\mathbf{R}^{n})\) ，且对每个 \(\beta \in \mathbf{N}^{n}\) （满足 \(|\beta| \leqslant 1\) ），广义函数 \(\partial^{\beta} f\) 属于 \(\mathrm{H}^{k-1}(\mathbf{R}^{n})\) 。由归纳假设，这等价于 \(f \in \mathrm{L}^{2}(\mathbf{R}^{n})\) ，且对每个 \((1 + \mathrm{N}^{k-1})\mathcal{F}(\partial^{\beta} f) \in \mathrm{L}^{2}(\mathbf{R}^{n})\) （满足 \(\beta \in \mathbf{N}^{n}\) ）有 \(|\beta| \leqslant 1\) 。由于 \(\mathcal{F}(\partial^{\beta} f) = (2i\pi X)^{|\beta|}\mathcal{F}(f)\) ，条件 \(f \in \mathrm{H}^{k}(\mathbf{R}^{n})\) 等价于说： \(\mathcal{F}f \in \mathrm{L}^{2}(\mathbf{R}^{n})\) ，且对每个 \((1 + \mathrm{N}^{k-1})\mathrm{X}^{\beta}\mathcal{F}f \in \mathrm{L}^{2}(\mathbf{R}^{n})\) （满足 \(\beta \in \mathbf{N}^{n}\) ）有 \(|\beta| \leqslant 1\) 。

不等式

\[
1 + \mathrm{N}^{k}\leqslant 1 + \mathrm{N}^{k - 1}\sum_{\substack{\beta \in \mathbf{N}^{n}\\ |\beta |\leqslant 1}}|\mathrm{X}^{\beta}|\leqslant 1 + n^{1 / 2}\mathrm{N}^{k}\leqslant (1 + n^{1 / 2})(1 + \mathrm{N}^{k})
\]

于是蕴含 \(f \in \mathrm{H}^k(\mathbf{R}^n)\) 当且仅当 \((1 + \mathrm{N}^k)\mathcal{F} \in \mathrm{L}^2(\mathbf{R}^n)\) 。

\section{部分算子}

\subsection{部分算子}

在本节中， K 是交换域。

我们回顾（E, II, §3, p. 9–10）：\(\text{graph}^{(2)}\) 是这样一个集合，其元素都是有序偶。设 A 与 B 是集合， A 与 B 之间的一个 \(\text{correspondence}\) 是一个三元组 \((\Gamma, A, B)\) ，其中 \(\Gamma\) 是包含在 A × B 中的一个图像；它的定义集（也称为 \(\text{domain}\) ）是 \(\text{pr}_1(\Gamma)\) ，而它的值集是 \(\text{pr}_2(\Gamma)\) 。一个对应是一个 \(\text{function}\) （E, II, p. 13，定义 9），如果它的图像是泛函的，且它的出发集与它的定义域重合。泛函图像的每个子集都是泛函图像。

定义 1. --- 设 E 与 F 是 K 上的向量空间。从 E 到 F 的部分算子 u 是 E 与 F 之间的一个对应 (Γ, E, F) ，它满足下述条件：

\begin{enumerate}
\def\labelenumi{(\roman{enumi})}
\item
  图像 \(\Gamma\) 是 \(\mathrm{E} \times \mathrm{F}\) 的一个向量子空间；
\item
  图像 \(\Gamma\) 是泛函的。
\end{enumerate}

若 E = F ，就说 u 是 E 上的部分算子。

设 u 是从 E 到 F 的部分算子。对应 u 的图像 \(\Gamma\) 称为部分算子 u 的图像，也记作 \(\Gamma_{u}\) 。我们用 \(\mathcal{P}(\mathrm{E};\mathrm{F})\) 表示从 E 到 F 的部分算子所成的集；我们简记 \(\mathcal{P}(\mathrm{E}) = \mathcal{P}(\mathrm{E};\mathrm{E})\) 。

给出一个从 E 到 F 的部分算子，相当于给出 E 的一个向量子空间 D 与一个从 D 到 F 的线性映射 u ；相应的部分算子是对应 \((\Gamma, \mathrm{E}, \mathrm{F})\) ，其中 \(\Gamma \subset D \times F\) 是 u 的图像。

部分算子 \(u\) 的定义域简称为 \(u\) 的定义域，并记作 \(\operatorname{dom}(u)\) 。

每个从 E 到 F 的线性映射 u 都是从 E 到 F 的部分算子。

若 D ⊂ E 是向量子空间，我们将用 1 \(_{D}\) 表示以 D 为定义域、在 D 上为恒等映射的部分算子，即对应 (Δ \(_{D}\) , E, E) ，其中 Δ \(_{D}\) 是 D × D 的对角线（E, II, p. 13，定义 8）。我们将用 \(0_{\mathrm{D}}\) 表示以 D 为定义域、在 D 上为零的部分算子，即对应 \((\mathrm{D} \times \{0\}, \mathrm{E}, \mathrm{F})\) 。

部分算子 \(u = (\Gamma, \mathrm{E}, \mathrm{F})\) 与 \(u' = (\Gamma', \mathrm{E}, \mathrm{F})\) （都是从 E 到 F 的）相等，当且仅当 \(\operatorname{dom}(u) = \operatorname{dom}(u')\) ，且从 \(u\) 到 F 中的线性映射 \(u'\) 与 \(\operatorname{dom}(u)\) 重合。

依照 E, II, §3 ，定义下列概念：

\begin{enumerate}
\def\labelenumi{(\roman{enumi})}
\tightlist
\item
  u 下 E 的子集 A 的像是 F 的子集 \(u(A \cap D)\) ；简记为 \(u(A)\) 。 u 下 F 的子集 B 的原像是 D 的子集 \(\overline{u}(B)\) 。
\end{enumerate}

若 A （相应地， B ）是 E （相应地， F ）的向量子空间，则它在 u 下的像（相应地，原像）是 F （相应地， E ）的向量子空间。

 \(u\) 的像是 \(u(\mathrm{D})\) 的向量子空间 \(\mathrm{F}\) ，也记作 \(\operatorname{Im}(u)\) 。称 \(u\) 为满射部分算子，如果 \(\operatorname{Im}(u) = \mathrm{F}\) 。 \(u\) 的核是 \(u^{-1}(\{0\})\) 的向量子空间 \(\mathrm{E}\) ，也记作 \(\operatorname{Ker}(u)\) 。 \(u\) 的核化为 0 ，当且仅当从 \(u\) （ \(\operatorname{dom}(u)\) 到 \(\mathrm{F}\) 中的）线性映射是单射的。这时称 \(u\) 是单射的。若 \(u\) 既单射又满射，就称它是双射的。

\begin{enumerate}
\def\labelenumi{(\roman{enumi})}
\setcounter{enumi}{1}
\item
  若 E 、 F 与 G 是 K 上的向量空间，且 \(u = (\Gamma, \mathrm{E}, \mathrm{F})\) ， \(v = (\Gamma', \mathrm{F}, \mathrm{G})\) 分别是从 E 到 F 与从 F 到 G 的部分算子，则复合对应 \(v \circ u = (\Gamma' \circ \Gamma, \mathrm{E}, \mathrm{G})\) 是从 E 到 G 中的部分算子。它的定义域是 \(u^{-1}(\mathrm{dom}(v))\) 。若 H 是 K 上的向量空间，且 \(w = (\Gamma'', \mathrm{G}, \mathrm{H})\) 是从 G 到 H 中的部分算子，则有 \(w \circ (v \circ u) = (w \circ v) \circ u\) 。有时把 \(v \circ u\) 简写作 vu 。
\item
  特别地，对每个从 E 到 F 中的部分算子 \(u\) 以及每个 \(a \in \mathrm{K}\) ，部分算子 \(au = (a1_{\mathrm{F}}) \circ u\) 与 \(ua = u \circ (a1_{\mathrm{E}})\) 有定义。若 \(a \neq 0\) 或者 \(u\) 的定义域等于 E ，则二者相等；我们有 \(u0 = 0_{\mathrm{E}}\) 与 \(0u = 0_{\mathrm{dom}(u)}\) 。
\end{enumerate}

设 E 是向量空间。由前所述，集合 \(\mathcal{P}(\mathrm{E})\) （赋有由 \((u,v)\mapsto u\circ v\) 所定义的复合律）是一个结合的含幺群胚（A, I, p. 4，定义 5 与 A, I, p. 12，定义 2），其单位元为 \(1_{\mathrm{E}}\) 。对每个 \(n\in \mathbf{N}\) ，我们用 \(u^n\) 表示复合 \({}^n u\) （A, I, p. 13）。

此外，我们定义下列概念：

\begin{enumerate}
\def\labelenumi{(\roman{enumi})}
\item
  若 \(u = (\Gamma, \mathrm{E}, \mathrm{F})\) 是从 E 到 F 的部分算子， G 是 E 的向量子空间，则 \(u\) 到 G 上的约化是从 E 到 F 的部分算子（ \(\Gamma \cap (\mathrm{G} \times \mathrm{F}), \mathrm{E}, \mathrm{F}\) ）。它的定义域是 \(\operatorname{dom}(u) \cap \mathrm{G}\) ；当不致与 \(u|_{\mathrm{G}}\) 到子空间 G 上的限制相混淆时，有时把它记作 \(u\) 。
\item
  设 \(v\) 是从 F 到 E 的单射部分算子。这时 \(v^{-1} = (\Gamma^{-1}, \mathrm{E}, \mathrm{F})\) 的逆对应 \(v\) 是一个部分算子，满足 \(\operatorname{dom}(v^{-1}) = \operatorname{Im}(v)\) 。称 \(v^{-1}\) 为 \(v\) 的逆部分算子。我们有等式 \(v \circ v^{-1} = 1_{\operatorname{dom}(v^{-1})}\) （在 \(\mathcal{P}(\mathrm{E})\) 中）与 \(v^{-1} \circ v = 1_{\operatorname{dom}(v)}\) （在 \(\mathcal{P}(\mathrm{F})\) 中）。部分算子 \(v^{-1}\) 是单射的，且有 \((v^{-1})^{-1} = v\) 。
\item
  设 E 、 F 与 G 是向量空间。设 u （相应地， v ）是从 E 到 F （相应地，从 F 到 G ）的单射部分算子。那么部分算子 \(v \circ u\) 是单射的，且 \((v \circ u)^{-1} = u^{-1} \circ v^{-1}\) 。
\item
  若 u 与 v 是从 E 到 F 的部分算子，我们说 v 是 u 的扩张，并记 \(u \subset v\) ，如果 u 的图像包含在 v 的图像中。这蕴含 \(\text{dom}(u) \subset \text{dom}(v)\) ，且 u 是 v 到 \(\text{dom}(u)\) 上的约化。关系 “ \(u \subset v\) ” 是 \(\mathcal{P}(\text{E};\text{F})\) 上的一个序关系。例如，对所有 \(au \subset ua\) 与每个 \(a \in K\) ，有 \(u \in \mathcal{P}(\text{E};\text{F})\) 。
\item
  设 E 是 K 上的向量空间， \((\mathrm{F}_{i})_{i\in\mathrm{I}}\) 是 K 上向量空间的一个族。对 \(i\in I\) ，设 \(u_{i}\) 是从 E 到 \(F_{i}\) 的一个部分算子。诸 \(u_{i}\) 的乘积部分算子是从 E 到诸空间 \(F_{i}\) 的乘积向量空间中的部分算子，其定义域是诸 \(\mathrm{dom}(u_{i})\) 的交 D ，它把 \(x\in D\) 对应到族 \((u_{i}(x))_{i\in\mathrm{I}}\) 。把它记作 \((u_{i})_{i\in\mathrm{I}}\) 。
\item
  设 A: F×F → F 是线性映射 \((x,y)\mapsto x+y\) 。设 u 与 v 是从 E 到 F 的部分算子。和 \(u+v\) 是从 E 到 F 的部分算子 A ∘ \((u,v)\) 。它的定义域是 dom \((u)\cap\text{dom}(v)\) 。对 \(\mathcal{P}(\text{E};\text{F})\) 中的 u 、 v 、 w ，有 \((u+v)+w=u+(v+w)\) 。
\end{enumerate}

设 G 是 K 上的向量空间。对 \(\mathcal{P}(\mathrm{E};\mathrm{F})\) 中所有 u 与 v 以及每个 \(w \in \mathcal{P}(\mathrm{F};\mathrm{G})\) ，有 \(w \circ u + w \circ v \subset w \circ (u + v)\) 。一般而言，这一公式中没有等号（IV, p. 344 的习题 1），但当 w 的定义域等于 F 时就是这种情形。对 \(w \in \mathcal{P}(\mathrm{G};\mathrm{E})\) ，有 \(u \circ w + v \circ w = (u + v) \circ w\) 。

\begin{enumerate}
\def\labelenumi{(\roman{enumi})}
\setcounter{enumi}{6}
\tightlist
\item
  设 L 是域 K 的一个扩张。设 E 与 F 是 K 上的向量空间， \(\mathrm{E}_{(\mathrm{L})} = \mathrm{L} \otimes_{\mathrm{K}} \mathrm{E}\) ， \(\mathrm{F}_{(\mathrm{L})} = \mathrm{L} \otimes_{\mathrm{K}} \mathrm{F}\) 是由 K 到 L 的纯量扩张得到的 L-向量空间（A, II, p. 82）。对每个从 E 到 F 的部分算子 u ，我们用 \(u_{(L)}\) 表示从 \(\mathrm{E}_{(\mathrm{L})}\) 到 \(\mathrm{F}_{(\mathrm{L})}\) 中的、其图像为 \(L \otimes_{K} \Gamma_{u}\) （ \(\mathrm{E}_{(\mathrm{L})} \times \mathrm{F}_{(\mathrm{L})}\) 的向量子空间）的部分算子；它的定义域是 \(L \otimes_{K} \operatorname{dom}(u)\) ，且它在这一定义域上与把 \(1 \otimes x\) 映到 \(1 \otimes u(x)\) （对所有 \(x \in \operatorname{dom}(u)\) ）的唯一线性映射重合。
\end{enumerate}

设 v 是从 E 到 F 的部分算子。我们有 \(u \subset v\) 当且仅当 \(u_{(L)} \subset v_{(L)}\) 。

\begin{enumerate}
\def\labelenumi{(\roman{enumi})}
\setcounter{enumi}{7}
\tightlist
\item
  设 \(E_{1}\) ， \(F_{1}\) ， \(E_{2}\) ， \(F_{2}\) 都是 K-向量空间。设 u （相应地， v ）是从 \(E_{1}\) 到 \(F_{1}\) （相应地，从 \(E_{2}\) 到 \(F_{2}\) ）的部分算子。我们用 \(u \otimes v\) 表示从 \(E_{1} \otimes F_{1}\) 到 \(E_{2} \otimes F_{2}\) 中的、定义域为 dom(u) \(\otimes\) dom(v) 且满足 \((u \otimes v)(x \otimes y) = u(x) \otimes v(y)\) （对所有 \((x, y) \in E_{1} \times E_{2}\) ）的部分算子。
\end{enumerate}

\subsection{闭算子、可闭算子与稠定算子}

在本节中， K 表示一个交换拓扑域（TG, III, p. 54）。

定义 2. --- 设 E 与 F 是 K 上的拓扑向量空间（EVT, I, p. 1，定义 1）。设 \(u \in \mathcal{P}(\mathrm{E};\mathrm{F})\) 是从 E 到 F 的部分算子。

若 u 的定义域在 E 中稠密，就称 u 是稠定算子。

若 u 的图像在拓扑向量空间 E × F 中是闭的，就称 u 是闭算子。若 u 有一个闭扩张，就称 u 是可闭的。

设 E 、 F 与 G 是 K 上的拓扑向量空间。稠定算子 \(u \in \mathcal{P}(\mathrm{E};\mathrm{F})\) 的每个扩张都是稠定的。此外，若 \(v \in \mathcal{L}(\mathrm{E};\mathrm{F})\) ，则 \(u + v\) 是稠定算子。若 \(v: F \to G\) （相应地， \(w: G \to E\) ）是拓扑向量空间的同构，则 \(v \circ u\) （相应地， \(u \circ w\) ）是稠定算子。

例. --- 设 E 与 F 是 K 上的拓扑向量空间。

\begin{enumerate}
\def\labelenumi{\arabic{enumi})}
\item
  设空间 F 是豪斯多夫的。设 u 是从 E 到 F 中的线性映射。若 u 连续，则部分算子 u 是闭的（TG, I, p. 53，推论 2）。再假设 K 是非离散赋值域，且 E 与 F 是 K 上的可度量化的完备拓扑向量空间。由闭图像定理（EVT, I, p. 19，推论 5），这时由 u 所定义的部分算子是闭的，当且仅当 u 连续。
\item
  设空间 E 是豪斯多夫的。设 v 是从 F 到 E 中的单射连续线性映射。这时部分算子 \(v^{-1} \in \mathcal{P}(\mathrm{E};\mathrm{F})\) （参见 IV, p. 226）是闭的，因为它的图像是 v 的图像在从 \(F \times E\) 到 \(E \times F\) 的、由 \((y,x) \mapsto (x,y)\) 所定义的拓扑向量空间同构下的像，而 v 的图像是闭的。
\end{enumerate}

命题 1. --- 设 E 与 F 是 K 上的拓扑向量空间。从 E 到 F 中的部分算子 u 是可闭的，当且仅当 u 的图像 \(\Gamma_{u}\) 在 E × F 中的闭包是一个泛函图像。这时存在唯一的从 E 到 F 的部分算子 v ，其图像为 \(\overline{\Gamma}_{u}\) ，且它是 u 的最小闭扩张。

若 \(u\) 的图像在 \(\mathrm{E} \times \mathrm{F}\) 中的闭包是泛函图像，则它就是某个部分算子的图像，而后者是 \(u\) 的一个闭扩张，因而它是可闭的。反之，设 \(u \subset w\) ，其中 \(w\) 是闭的。闭包 \(\overline{\Gamma}_u\) （即 \(u\) 的图像在 \(\mathrm{E} \times \mathrm{F}\) 中的闭包）包含在 \(\Gamma_w\) 中，因而 \(\overline{\Gamma}_u\) 是泛函图像。

最后的断言由下述事实得出：若 w 是 u 的一个闭扩张，则 w 的图像包含 \(\overline{\Gamma}_{u}\) 。

定义 3. --- 设 E 与 F 是 K 上的拓扑向量空间。设 u 是从 E 到 F 中的可闭算子。图像为 \(\overline{\Gamma}_{u}\) 的闭算子称为 \(u\) 的闭包，记作 \(\overline{u}\) 。

注. --- 设 E 与 F 是 K 上的拓扑向量空间。设 u 是从 E 到 F 的可闭算子。u 的闭包的定义域包含在 u 的定义域在 E 中的闭包内。一般而言，二者并不相等（IV, p. 344 的习题 1, b)）。

若 \(u \in \mathcal{P}(\mathrm{E};\mathrm{F})\) 是可闭的且 \(\operatorname{dom}(u) = \operatorname{E}\) ，则 \(u = \overline{u}\) 是闭的，因为这时 \(\operatorname{dom}(\overline{u}) = \operatorname{dom}(u)\) 。

命题 2. --- 设 K = R ，并设 E 与 F 是 R 上的拓扑向量空间。设 u 是从 E 到 F 中的部分算子。则 u 是稠定的（相应地，是闭的、是可闭的），当且仅当 \(u_{(\mathbf{C})}\) 是稠定的（相应地，是闭的、是可闭的）。

设 \(u\) 的定义域在 E 中稠密。 \(\mathrm{E}_{(\mathbf{C})}\) 中 0 的每个邻域都包含形如 \(\mathrm{V} + i\mathrm{V}\) 的一个邻域（EVT, II, p. 65），其中 V 是 E 中 0 的一个邻域，因而它包含 \(u_{(\mathbf{C})}\) 的定义域中的一个元素；所以这一部分算子是稠定的。反过来也成立，因为从 \(\mathrm{E}_{(\mathbf{C})}\) 到 E 中的、把 \(x + iy\) 映到 \(x\) （对所有 \((x,y) \in \mathrm{E} \times \mathrm{E}\) ）的映射是连续且满射的。

我们将 \(u_{(\mathbf{C})}\) 的图像等同于 \(u\) 的图像的复化拓扑向量空间。于是，它在 \(\mathrm{E}_{(\mathbf{C})} \times \mathrm{E}_{(\mathbf{C})}\) 中是闭的，如果 \(u\) 的图像在 \(\mathrm{E} \times \mathrm{E}\) 中是闭的。反之，我们有 \(\Gamma_u = \Gamma_{u_{(\mathbf{C})}} \cap (\mathrm{E} \times \mathrm{E})\) （在 \(\mathrm{E}_{(\mathbf{C})} \times \mathrm{E}_{(\mathbf{C})}\) 中）；由于 \(\mathrm{E} \times \mathrm{E}\) 在 \(\mathrm{E}_{(\mathbf{C})} \times \mathrm{E}_{(\mathbf{C})}\) 中是闭的，故部分算子 \(u\) 在 \(u_{(\mathbf{C})}\) 闭时是闭的。

从 E 到 F 中的部分算子 \(v\) 是 \(u\) 的扩张，当且仅当 \(v_{(\mathbf{C})}\) 是 \(u_{(\mathbf{C})}\) 的扩张；因此，部分算子 \(u_{(\mathbf{C})}\) 在 \(u\) 可闭时是可闭的。反之，若 \(u_{(\mathbf{C})}\) 可闭，则部分算子 \(u\) 也可闭，因为 \(\overline{\Gamma_u} = \overline{\Gamma_{u_{(\mathbf{C})}}} \cap (\mathrm{E} \times \mathrm{E})\) 这时是泛函图像（命题 1）。

引理 1. — 设 E 、 F 与 G 是 K 上的拓扑向量空间。设 u 是从 E 到 F 的闭算子。

\begin{enumerate}
\def\labelenumi{\alph{enumi})}
\item
  对每个 \(v \in \mathcal{L}(\mathrm{E};\mathrm{F})\) ，部分算子 \(u + v\) 是闭的。
\item
  对每个 \(v \in \mathcal{L}(\mathrm{G};\mathrm{E})\) ，部分算子 \(u \circ v\) 是闭的。
\end{enumerate}

我们来证明 \(a\) )。设 \(\gamma\) 是映射 \((x,y)\mapsto (x,y - v(x))\) ，它从 \(\mathrm{E}\times \mathrm{F}\) 到其自身，并且是连续的。对每个 \((x,y)\in \mathrm{E}\times \mathrm{F}\) ，我们有 \(\gamma (x,y)\in \Gamma_u\) 当且仅当 \(x\in \operatorname {dom}(u)\) 且 \(y = u(x) + v(x)\) ，亦即 \(\overline{\gamma}^{1}(\Gamma_{u}) = \Gamma_{u + v}\) 。断言得证。

我们来证明 \(b\) )。映射 \(\eta = (v, 1_{\mathrm{F}})\) 从 \(\mathbf{G} \times \mathbf{F}\) 到 \(\mathbf{E} \times \mathbf{F}\) 中是连续的；对每个 \((z, y) \in \mathbf{G} \times \mathbf{F}\) ，我们有 \(\eta(z, y) = (v(z), y)\) ，从而 \(\overline{\eta}(\Gamma_u) = \Gamma_{u \circ v}\) 。断言得证。

设 E 与 F 是 K 上的拓扑向量空间，且 F 是豪斯多夫的。设 \(a \in \mathrm{K}\) 。若 \(u \in \mathcal{P}(\mathrm{E};\mathrm{F})\) 是可闭的，则 \(au\) 也是可闭的。若 \(a \neq 0\) ，我们有 \(\overline{au} = a\overline{u}\) ，且 \(u\) 是闭的当且仅当 \(au\) 是闭的。若 \(a = 0\) ，则 \(au\) 的闭包是 \(0_{\overline{\operatorname{dom}(u)}}\) ，而 \(a\overline{u}\) 等于 \(0_{\operatorname{dom}(\overline{u})}\) ；因此可能出现这样的情形： \(u\) 是闭的，而 \(au\) 不是闭的。

命题 3. --- 设 E 与 F 是 K 上的拓扑向量空间，其中 F 是豪斯多夫的。设 u 是从 E 到 F 的闭部分算子。u 的核是 E 的一个闭向量子空间。

事实上，u 的核是 E × F 的闭子空间 \(\Gamma_{u} \cap (\mathrm{E} \times \{0\})\) 在从 E 到 E × F 中的连续线性映射 \(x \mapsto (x, 0)\) 下的原像。

在本节余下部分，我们假设 K 是非离散赋值域。

定义 4. --- 设 E 与 F 是 K 上的赋范空间，u 是从 E 到 F 的部分算子。对 dom(u) 中的 x ，我们记

\[
\| x \| _ {u} = (\| x \| _ {\mathrm{E}} ^ {2} + \| u (x) \| _ {\mathrm{F}} ^ {2}) ^ {1 / 2}.
\]

如此定义的映射 \(x \mapsto \|x\|_{u}\) 是 \(\text{dom}(u)\) 上的一个范数。我们把这样得到的赋范空间记作 \(E_{u}\) 。

注. --- 设 E 与 F 是 K 上的赋范空间，u 是从 E 到 F 中的部分算子。

\begin{enumerate}
\def\labelenumi{\arabic{enumi})}
\item
  \(E_{u}\) 到 E 中的典则单射是连续的，因为我们有 \(\|x\|\leqslant\|x\|_{u}\) （对所有 \(x\in E\) ）。特别地，\(\operatorname{dom}(u)\) 的在 E 中闭的每个子空间在 \(E_{u}\) 中也是闭的。
\item
  若 E 与 F 是希尔伯特空间，则空间 \(E_{u}\) 是准希尔伯特空间，因为 \(E_{u}\) 上的范数来自下述正埃尔米特形式
\end{enumerate}

\[
(x, y) \mapsto (x \mid y) _ {u} = \langle x \mid y \rangle + \langle u (x) \mid u (y) \rangle
\]

（它定义在 dom(u) 上）。

命题 4. --- 设 E 与 F 是巴拿赫空间（相应地，希尔伯特空间），u 是从 E 到 F 中的部分算子。则 u 是闭的，当且仅当赋范空间 \(E_{u}\) 是巴拿赫空间（相应地，希尔伯特空间）。

只须处理巴拿赫空间的情形。 \(E_{u}\) 的范数，是经由双射线性映射 \((x,y)\mapsto x\) （它把 \(\Gamma_{u}\) 映到 dom(u) 中）作结构迁移，从下述范数得到的：即限制到子空间 \(\Gamma_{u}\) （赋以范数 \((x,y)\mapsto(\|x\|_{\mathrm{E}}^{2}+\|y\|_{\mathrm{F}}^{2})^{1/2}\) ，它由巴拿赫空间 \(E\oplus F\) 上的范数限制而得）上所得到的范数。因此，空间 \(E_{u}\) 是巴拿赫空间，当且仅当子空间 \(\Gamma_{u}\) （在 \(E\oplus F\) 中）是闭的。

若 \(u\) 与 \(v\) 分别是从 E 到 F 与从 F 到 G 的部分算子，且若 \(u\) 是闭的，则部分算子 \(v \circ u\) 一般说来不是闭的，即使 \(v\) 是连续的（IV, p. 344 的习题 1, c)）。尽管如此，我们有如下的充分条件：

引理 2. --- 设 E 、 F 与 G 是 K 上的赋范空间，其中 F 是巴拿赫空间。设 u 是从 E 到 F 的闭部分算子，并设 \(v \in \mathcal{L}(\mathrm{F};\mathrm{G})\) 。如果存在 \(\mathrm{C} \in \mathbf{R}_{+}\) 使得

\[
\| u (x) \| \leqslant C (\| x \| + \| (v \circ u) (x) \|)
\]

对所有 \(x \in \text{dom}(v \circ u) = \text{dom}(u)\) 成立；也就是说，若从 \(x \mapsto u(x)\) 到 F 中的线性映射 \(E_{v \circ u}\) 是连续的，则 \(v \circ u\) 是闭的。

记 \(w = v \circ u\) 。设 \((x_n, w(x_n))_{n \in \mathbf{N}}\) 是 \(w\) 的图像中的一个序列，它在 \(E \times G\) 中收敛。设 \(x\) 是序列 \((x_n)\) 的极限。假设表明，序列 \((u(x_n))_{n \in \mathbf{N}}\) 这时是 \(F\) 中的一个柯西序列；它收敛到元素 \(y\) ，该元素属于 \(F\) 。由于 \(u\) 是闭的，我们有 \(x \in \text{dom}(u)\) 与 \(y = u(x)\) 。于是 \(w(x_n) = v(u(x_n))\) 趋于 \(v(y) = v(u(x))\) ，因为 \(v\) 连续；因此 \(w\) 的图像是闭的。

设 u 是巴拿赫空间 E 上的闭部分算子， F 是 \(\text{dom}(u)\) 的一个子空间。若 F 在巴拿赫空间 \(E_{u}\) 中稠密，则 u 到 F 上的约化是可闭的，且其闭包等于 u 。这时称 F 为 u 的一个核。

\subsection{部分算子的例子}

在本节中， K = R 或 C 。

例. --- 1) 设 X 是局部紧拓扑空间， \(\mu\) 是 X 上的一个正测度。我们固定元素 \(p_1\) 与 \(p_2\) ，它们属于 \([1, +\infty[\) 。

设 \(g\) 是 \(\mu\)-可测的函数，定义在 \(X\) 上，取值于 \(K\) 中。用 \(D\) 表示 \(\mathcal{L}_{K}^{p_1}(X, \mu)\) 中满足如下条件的函数 \(f\) （属于 \(\mathcal{L}_{K}^{p_1}(X, \mu)\) ）所成的子空间： \(gf \in \mathcal{L}_{K}^{p_2}(X, \mu)\) 。由 \(D\) 到 \(\mathcal{L}_{K}^{p_2}(X, \mu)\) 中的、由 \(f \mapsto gf\) 所定义的线性映射确定了 \(\mathcal{L}_{K}^{p_1}(X, \mu)\) 到 \(\mathcal{L}_{K}^{p_2}(X, \mu)\) 中的一个部分算子，我们把它记作 \(m_g\) 。

\(\mu\)-可忽略函数在 \(\mathcal{L}_{\mathrm{K}}^{p_1}(\mathrm{X},\mu)\) 中所成的向量子空间包含在 D 中；并且， \(m_g\) 把 \(\mu\)-可忽略函数映为仍为 \(\mu\)-可忽略的函数。我们将用 \(\tilde{m}_g\) 表示从 \(\mathrm{L}_{\mathrm{K}}^{p_1}(\mathrm{X},\mu)\) 到 \(\mathrm{L}_{\mathrm{K}}^{p_2}(\mathrm{X},\mu)\) 中的、由 \(m_g\) 通过过渡到商诱导出的部分算子。它称为以 \(g\) 为乘子的乘法算子（从 \(\mathrm{L}_{\mathrm{K}}^{p_1}(\mathrm{X},\mu)\) 到 \(\mathrm{L}_{\mathrm{K}}^{p_2}(\mathrm{X},\mu)\) 中）。凡是 \(g_1\) 与 \(g_2\) 局部 \(\mu\)-几乎处处相等的函数，都定义同一个乘法算子。

命题 5. — 乘法算子 \(\widetilde{m}_{g}\) （从 \(\mathrm{L}_{\mathrm{K}}^{p_{1}}(\mathrm{X},\mu)\) 到 \(\mathrm{L}_{\mathrm{K}}^{p_{2}}(\mathrm{X},\mu)\) ）是稠定的闭算子。

我们先证明部分算子 \(\widetilde{m}_g\) 是闭的。设 \((f_n, h_n)_{n \in \mathbf{N}}\) 是 \(\mathcal{L}_{\mathrm{K}}^{p_1}(\mathrm{X}, \mu) \times \mathcal{L}_{\mathrm{K}}^{p_2}(\mathrm{X}, \mu)\) 中的一个序列，使得 \((\widetilde{f}_n, \widetilde{h}_n)\) 的类所作成的序列（即 \(f_n\) 与 \(h_n\) 的类）属于 \(\widetilde{m}_g\) 的图像，且当 \(\mathrm{L}_{\mathrm{K}}^{p_1}(\mathrm{X}, \mu) \times \mathrm{L}_{\mathrm{K}}^{p_2}(\mathrm{X}, \mu)\) 趋于无穷时在 \(n\) 中收敛。设 \((f, h)\) 是 \(\mathcal{L}_{\mathrm{K}}^{p_1}(\mathrm{X}, \mu) \times \mathcal{L}_{\mathrm{K}}^{p_2}(\mathrm{X}, \mu)\) 中的一对元素，使得它们的类所成的对 \((\widetilde{f}, \widetilde{h})\) 是 \((\widetilde{f}_n, \widetilde{h}_n)\) 的极限。

存在序列 \((f_{n_{k}})_{k\in\mathbf{N}}\) （从序列 \((f_{n})_{n}\) 中抽取的）使得 \(f_{n_{k}}(x)\) 收敛到 \(f(x)\) ，对 \(\mu\)-几乎所有的 x 成立（INT, IV, p. 131, § 3, no. 4，定理 3）。由此推出， \(h_{n_{k}}(x)=g(x)f_{n_{k}}(x)\) \(\mu\)-几乎处处收敛于 \(g(x)f(x)\) 。此外，序列 \((h_{n_{k}})\) 在空间 \(\mathcal{L}_{\mathrm{K}}^{p_{2}}(\mathrm{X},\mu)\) 中收敛于 h 。因此，函数 h 与 gf \(\mu\)-几乎处处相等（loc. cit.）。于是 \(\widetilde{f}\) 属于 \(\widetilde{m}_{g}\) 的定义域，且 \(\widetilde{h}=\widetilde{m}_{g}(\widetilde{f})\) 。这就证明了 \(\widetilde{m}_{g}\) 是闭的。

我们来证明 \(\widetilde{m}_g\) 的定义域在 \(\mathrm{L}_{\mathrm{K}}^{p_1}(\mathrm{X},\mu)\) 中稠密。只须验证：函数 \(f\in \mathcal{K}(\mathrm{X};\mathrm{K})\) 的类属于 \(\widetilde{m}_g\) 的定义域在 \(\mathrm{L}_{\mathrm{K}}^{p_1}(\mathrm{X},\mu)\) 中的闭包。设 \(f\in \mathcal{K}(\mathrm{X};\mathrm{K})\) ，并用 \(\widetilde{f}\) 表示它在 \(\mathrm{L}_{\mathrm{K}}^{p_1}(\mathrm{X},\mu)\) 中的类。对每个整数 \(n\in \mathbf{N}\) ，用 \(\varphi_{n}\) 表示由元素 \(x\in \mathrm{X}\) （满足 \(|g(x)|\leqslant n\) ）所成集合的特征函数，并令 \(f_{n} = f\varphi_{n}\) 。这时我们有 \(|gf_n|\leqslant n|f|\) ，而它属于 \(\mathcal{L}_{\mathrm{K}}^{p_2}(\mathrm{X},\mu)\) ，因此 \(f_{n}\) 属于 \(\widetilde{m}_g\) 的定义域。对 X 的每个元素 \(x\) ，序列 \((f_n(x))_{n\in \mathbf{N}}\) 收敛于 \(f(x)\) （当 \(n\to +\infty\) 时）；此外，我们有 \(|f_n|\leqslant |f|\) ，它属于 \(\mathcal{L}_{\mathrm{K}}^{p_1}(\mathrm{X},\mu)\) 。由勒贝格定理（INT, IV, p. 137, § 3, n°7，定理 6）， \(f_{n}\) 的类的序列在 \(\widetilde{f}\) 中收敛于 \(\mathrm{L}_{\mathrm{K}}^{p_1}(\mathrm{X},\mu)\) 。因此， \(f\) 的类属于 \(\widetilde{m}_g\) 的定义域的闭包。

在下述命题中，我们假设 \(p_1 = p_2 = 2\) 。

命题 6. --- a) 设 \(g'\) 是 X 上的 \(\mu\)-可测函数，满足 \(|g|\leqslant |g'|\) 。我们有 \(\mathrm{dom}(\widetilde{m}_{g'})\subset\mathrm{dom}(\widetilde{m}_{g})\) ，且 \(\mathrm{dom}(\widetilde{m}_{g'})\) 是部分算子 \(\widetilde{m}_{g}\) 的一个核；

\begin{enumerate}
\def\labelenumi{\alph{enumi})}
\setcounter{enumi}{1}
\tightlist
\item
  设 F 是 \(\mathcal{L}_{\mathrm{K}}^{2}(\mathrm{X},\mu)\) 的一个子空间，它与 \(\mathcal{K}(\mathrm{X};\mathrm{K})\) 的交在 \(\mathcal{K}(\mathrm{X};\mathrm{K})\) 中稠密，且它在 \(\mathrm{L}_{\mathrm{K}}^{2}(\mathrm{X},\mu)\) 中的像 G 包含在 \(\mathrm{dom}(\widetilde{m}_g)\) 中。若 \(|g|^2\) 局部 \(\mu\)-可积，则 \(\mathcal{K}(\mathrm{X};\mathrm{K})\) 包含在 \(\widetilde{m}_g\) 的定义域中，且 G 是 \(m_g\) 的一个核。
\end{enumerate}

我们来证明 \(a\) )。若 \(f \in \mathcal{L}_{\mathrm{K}}^{2}(\mathrm{X}, \mu)\) 属于 \(m_{g'}\) 的定义域（从而 \(fg' \in \mathcal{L}_{\mathrm{K}}^{2}(\mathrm{X}, \mu)\) ），则假设蕴含 \(fg \in \mathcal{L}_{\mathrm{K}}^{2}(\mathrm{X}, \mu)\) ，由此即得结论。

我们来证明 \(\mathrm{dom}(\widetilde{m}_{g'})\) 是 \(\widetilde{m}_g\) 的核，也就是说， \(\widetilde{m}_{g'}\) 的定义域在希尔伯特空间 \(\mathrm{E}_{\widetilde{m}_g}\) 中稠密。设 \(h \in \mathcal{L}_{\mathrm{K}}^2(\mathrm{X}, \mu)\) ，其类 \(\widetilde{h}\) 属于 \(\mathrm{E}_{\widetilde{m}_g}\) 且与 \(\mathrm{dom}(\widetilde{m}_{g'})\) 正交。这就意味着

\[
(\widetilde {h} \mid \widetilde {h} ^ {\prime}) _ {\widetilde {m} _ {g}} = \int_ {\mathrm{X}} \overline {{h}} h ^ {\prime} (1 + | g | ^ {2}) d \mu = 0
\]

它对每个这样的函数 \(h' \in \mathcal{L}_{\mathrm{K}}^{2}(\mathrm{X}, \mu)\) （其类 \(\widetilde{h}'\) 属于 \(\operatorname{dom}(\widetilde{m}_{g'})\) ）成立。

设 C 是 X 的紧子集， \(\varphi\) 是它的特征函数。设 \(n \in \mathbf{N}\) 。设 \(\varphi_n\) 是 \(\mu\)-可积集合 \(\mathrm{C}_n\) （由 C 中满足 \(x \in \mathrm{C}\) 的元素 \(|h(x)| \leqslant n\) 所成）的特征函数，并令 \(h_n' = \varphi_n h\) 。 \(h_n'\) 的类属于 \(\tilde{m}_{g'}\) 的定义域，因为 \(|g'h_n'| \leqslant n\varphi\) ；由此得出

\[
0 = \int_ {\mathrm{X}} \overline {{h}} h _ {n} ^ {\prime} (1 + | g | ^ {2}) d \mu = \int_ {\mathrm{X}} | h | ^ {2} \varphi_ {n} (1 + | g | ^ {2}) d \mu .
\]

这蕴含：对 \(\mu\)-几乎所有 \(x \in C_{n}\) ， h 为零；由于 n 任意，对 \(\mu\)-几乎所有 \(x \in C\) ， h 为零。最终推出 h \(\mu\) \(\mu\)-几乎处处为零，因为 C 任意且 h 是缓增的（INT, V, p. 9, § 1, n°3，推论）。

现在考虑断言 \(b\) )。由于 \(|g|^2\) 局部 \(\mu\)-可积，函数 \(fg\) 属于 \(\mathcal{L}_{\mathrm{K}}^{2}(\mathrm{X},\mu)\) （对 \(f\in\mathcal{K}(\mathrm{X};\mathrm{K})\) ），因此 \(\mathcal{K}(\mathrm{X};\mathrm{K})\) 包含在 \(\widetilde{m}_{g}\) 的定义域中。

设 \(h \in \mathcal{L}_{\mathrm{K}}^{2}(\mathrm{X}, \mu)\) ，其类 \(\widetilde{h}\) 属于 \(\mathrm{E}_{\widetilde{m}_{g}}\) 且与 G 正交。这时我们有

\[
0 = (\widetilde {h} \mid \widetilde {h} ^ {\prime}) _ {\widetilde {m} _ {g}} = \int_ {\mathrm{X}} h \overline {{h}} ^ {\prime} (1 + | g | ^ {2}) d \mu
\]

对所有 \(h' \in F\) ，其类为 \(h'\) 。鉴于关于 F 的假设，这意味着测度 \(h(1 + |g|^{2}) \cdot \mu\) 为零，因而 h \(\mu\)-几乎处处为零，因为 h 是缓增的。

设 \(p\) 是实数 \(\geqslant 1\) 。设 \(h\) 是 \(\mathcal{L}_{\mathrm{K}}^{\infty}(\mathrm{X},\mu)\) 中的元素。以 \(\widetilde{m}_h\) 为乘子的乘法算子 \(h\) 是 \(\mathrm{L}_{\mathrm{K}}^{p}(\mathrm{X},\mu)\) 的一个自同态。假设由 \(x\in X\) 中满足 \(h(x)=0\) 的 x 所成的集 Y 是局部 \(\mu\)-可忽略的。这时自同态 \(\widetilde{m}_h\) 是单射的（IV, p. 186 的引理 7）。用 \(h^{-1}\) 表示 X 上在 Y 上等于 0 、在 \(x \mapsto 1/h(x)\) 上等于 \(\mathrm{X} - \mathrm{Y}\) 的函数。逆部分算子 \(\widetilde{m}_{h}^{-1}\) 就是以 \(h^{-1}\) 为乘子的、从 \(\mathrm{L}_{\mathrm{K}}^{p}(\mathrm{X},\mu)\) 到 \(\mathrm{L}_{\mathrm{K}}^{p}(\mathrm{X},\mu)\) 中的乘法算子，即 \(\widetilde{m}_{h}^{-1} = \widetilde{m}_{h^{-1}}\) 。事实上， \(\widetilde{m}_{h}\) 的像是由形如 \(g \in \mathcal{L}_{\mathrm{K}}^{p}(\mathrm{X},\mu)\) （其中 \(g = hf\) ）的函数 \(f \in \mathcal{L}_{\mathrm{K}}^{p}(\mathrm{X},\mu)\) 的类所成的空间。这一条件等价于：对所有 \(g(x)/h(x) = f(x)\) 有 \(x \in \mathrm{X} - \mathrm{Y}\) ，且若 \(g(x) = 0\) 则 \(x \in \mathrm{Y}\) 。这蕴含： \(\widetilde{m}_{h}^{-1}\) 在 \(\mathrm{L}_{\mathrm{K}}^{p}(\mathrm{X},\mu)\) 中的定义域就是 \(\widetilde{m}_{h^{-1}}\) 的定义域，且等式 \(\widetilde{m}_{h}^{-1} = \widetilde{m}_{h^{-1}}\) 成立。

在下文中，我们有时把部分乘法算子简记为 \(m_{h}\) ，它以 h 为乘子，是从 \(\mathrm{L}_{\mathrm{K}}^{p_{1}}(\mathrm{X},\mu)\) 到 \(\mathrm{L}_{\mathrm{K}}^{p_{2}}(\mathrm{X},\mu)\) 中的。

\begin{enumerate}
\def\labelenumi{\arabic{enumi})}
\setcounter{enumi}{1}
\tightlist
\item
  设 E 是 K 上的希尔伯特空间， \(\mathrm{B} = (e_i)_{i \in \mathrm{I}}\) 是 E 的一个标准正交基。设 \((\lambda_i)_{i \in \mathrm{I}}\) 是 K 中元素的一个族。设 D 是 E 中由元素 \(x \in E\) （使得族 \((\lambda_i \langle e_i | x \rangle)_{i \in \mathrm{I}}\) 在 K 中平方可和）所成的向量子空间。空间 D 在 E 中稠密，因为它包含向量 \(e_i\) （对所有 \(i \in I\) ）。以 D 为定义域的部分算子 u 由下式给出
\end{enumerate}

\[
x \mapsto \sum_ {i \in I} \lambda_ {i} \langle e _ {i} | x \rangle e _ {i}
\]

称为基 B 中的部分对角算子，而 \((\lambda_{i})_{i\in I}\) 称为 u 的特征值族。

算子 \(u\) 是闭的。事实上，设 \((x_{n}, u(x_{n}))_{n \in \mathbf{N}}\) 是 \(u\) 的图像中在 \(\mathrm{E} \times \mathrm{E}\) 内收敛的一个元素序列，并设 \((x, y)\) 为其极限。这时我们有 \(\langle e_{i} | x_{n} \rangle \to \langle e_{i} | x \rangle\) （对所有 \(i \in \mathrm{I}\) ），并且

\[
\langle e _ {i} | u (x _ {n}) \rangle = \lambda_ {i} \langle e _ {i} | x _ {n} \rangle \rightarrow \langle e _ {i} | y \rangle
\]

对所有 \(i \in I\) 。于是， \(\lambda_{i} \langle e_{i} \mid x \rangle = \langle e_{i} \mid y \rangle\) 对所有 \(i \in I\) 成立；这就证明了 \(x \in D\) 且 \(u(x) = y\) ，即 \(u\) 是闭的。

这个例子实际上是前一例子应用于拓扑空间 X = I （赋以离散拓扑与计数测度 \(\mu\) ）的特殊情形，因为 E 通过映射 \(\ell^{2}(\mathrm{I}) = \mathrm{L}^{2}(\mathrm{I}, \mu)\) 等同于空间 \(x \mapsto (\langle e_{i} \mid x \rangle)_{i \in \mathrm{I}}\) （EVT, V, p. 23，推论 2），而 u 这时等同于乘法算子 \(m_{\lambda}\) ，其中 \(\lambda\) 是函数 \(i \mapsto \lambda_{i}\) 。

\begin{enumerate}
\def\labelenumi{\arabic{enumi})}
\setcounter{enumi}{2}
\tightlist
\item
  集合 \(N_{R} = N \cup \{\infty, \omega\}\) 赋有 VAR, R2, p. 10 中所述的全序，使得 \(n < \infty < \omega\) 对所有 \(n \in N\) 成立。设 \(r \in N_{R}\) 。设 \(n \in N\) ，并设 U 是 \(R^{n}\) 的一个开子集。设 \(k \in N\) 满足 \(k \leqslant r\) 。设 \((n_{\alpha})_{|\alpha| \leqslant k}\) 是 \(\mathcal{C}^{r}(U)\) 中元素的一个族，其中所考虑的重指标属于 \(N^{n}\) 。族 \((n_{\alpha})\) 在 U 上定义一个阶 \(\leqslant k\) 的纯量微分算子 D （参见 VAR, R2, 14.1.6, 14.1.4）。
\end{enumerate}

对每个满足 \(k \leqslant m \leqslant r\) 的整数 m ，微分算子 D 定义一个从 \(\mathrm{C}^{m}(\mathrm{U})\) 到 \(\mathrm{C}^{m-k}(\mathrm{U})\) 中的线性映射，它把 \(f \in \mathrm{C}^{m}(\mathrm{U})\) 映为

\[
\mathrm{D} (f) = \sum_ {| \alpha | \leqslant k} n _ {\alpha} \partial^ {\alpha} (f).
\]

同一公式定义了一个从 \(\mathcal{D}(\mathrm{U})\) 到 \(\mathcal{D}'(\mathrm{U})\) 中的连续线性映射（IV, p. 208 的定义 2）。

定义 5. --- 设 E 是 \(\mathcal{D}'(\mathrm{U})\) 的一个包含 \(\mathcal{D}(\mathrm{U})\) 的向量子空间。 E 上与 D 相伴的微分算子，指的是 E 上这样一个部分算子：它是以下述部分算子的扩张，该算子以 \(\mathcal{D}(\mathrm{U})\) 为定义域、由 \(\varphi \mapsto \mathrm{D}(\varphi)\) 所定义。

例如，设系数 \(n_{\alpha}\) 是 U 上的有界函数。设 \(\mu\) 是 U 上的勒贝格测度。若 p 是 \([1,+\infty[\) 中的一个元素，则可在 \(\mathrm{L}^{p}(\mathrm{U})\) 上定义一个与 D 相伴的微分算子，其定义域是索伯列夫空间 \(\mathrm{W}^{k,p}(\mathrm{U})\) （IV, p. 221 的第 14 节），因为这时我们有 \(n_{\alpha}\partial^{\alpha}(f)\in\mathrm{L}^{p}(\mathrm{U})\) （对所有 \(f\in\mathrm{W}^{k,p}(\mathrm{U})\) 与每个 \(|\alpha|\leqslant k\) ）。

\subsection{伴随}

在本节中， K 是域 R 或 C 之一， E 与 F 表示 K 上的希尔伯特空间。

设 \(u\) 是从 E 到 F 的稠定算子。用 \(D\) 表示 \(u\) 的定义域。对 \(y \in \mathrm{F}\) ，令 \(\lambda_y\) 为 D 上由 \(\lambda_y(x) = \langle y \mid u(x) \rangle\) （对所有 \(x \in \mathrm{D}\) ）所给的线性形式。我们用 \(\mathrm{D}^*\) 表示满足下述条件的向量 \(y \in \mathrm{F}\) 所成的集： \(\lambda_y\) 在 D 上连续。它是 F 的一个向量子空间。设 \(y \in \mathrm{D}^*\) ；由于 D 在 E 中稠密，线性形式 \(\lambda_y\) 唯一地扩张为 E 上的连续线性形式，我们仍把它记作 \(\lambda_y\) 。依照 EVT, V, p. 15，定理 3，存在唯一的元素 \(u^*(y)\) 属于 E ，使得 \(\lambda_y(x) = \langle u^*(y) \mid x \rangle\) 对所有 \(x \in \mathrm{E}\) 成立。映射 \(y \mapsto u^*(y)\) 是从 \(\mathrm{D}^*\) 到 E 中的线性映射。

定义 6. --- 以 \(D^*\) 为定义域、由 \(y \mapsto u^*(y)\) 所定义的从 F 到 E 的部分算子称为 \(u\) 的伴随，记作 \(u^*\) 。

注. --- 于是， \(y \in D^{*}\) 当且仅当存在 \(c \in R_{+}\) 使得 \(|\langle y | u(x) \rangle| \leqslant c \|x\|\) 对所有 \(x \in D\) 成立。这时元素 \(u^{*}(y)\) 由下述关系式刻画：

\[
\langle y \mid u (x) \rangle = \langle u ^ {*} (y) \mid x \rangle\tag{1}
\]

对所有 \(x \in \mathrm{D}\) 成立。这时我们有 \(|\langle y | u(x) \rangle| \leqslant \| u^{*}(y) \| \| x \|\) 对所有 \(x \in \mathrm{D}\) 成立。

在 \(u\) 是从 E 到 F 中的连续线性映射的情形，按前一定义意义的伴随与 EVT, V, p. 38，定义 1 中定义的伴随重合，因为这时 D* 等于 F 。

设 \(v \in \mathcal{P}(\mathrm{E};\mathrm{F})\) 满足 \(u \subset v\) 。这时我们有 \(v^{*} \subset u^{*}\) 。

设 \(v \in \mathcal{P}(\mathrm{E};\mathrm{F})\) 满足 \(u + v\) 稠定。这时部分算子 \(v\) 便是稠定的，且有 \(u^{*} + v^{*} \subset (u + v)^{*}\) 。一般而言等式不成立（IV, p. 347 的习题 9）。若 \(v \in \mathcal{L}(\mathrm{E};\mathrm{F})\) ，则 \(u + v\) 稠定，且 \((u + v)^{*} = u^{*} + v^{*}\) 。例如，若 \(\mathrm{F} = \mathrm{E}\) 且 \(v = \lambda 1_{\mathrm{E}}\) （其中 \(\lambda \in \mathrm{K}\) ），就是这种情形。

设 G 是 K 上的希尔伯特空间， \(v \in \mathcal{P}(\mathrm{F};\mathrm{G})\) 是稠定算子。若 \(v \circ u\) 稠定，则 \(u^{*} \circ v^{*} \subset (v \circ u)^{*}\) 。一般而言等式不成立（loc. cit.）。若 u （相应地， v ）是同构，则 \(v \circ u\) 稠定且有 \((v \circ u)^{*} = u^{*} \circ v^{*}\) 。例如，若 E = F （相应地， F = G ）且 \(u = \lambda 1_{E}\) （相应地， \(v = \lambda 1_{F}\) ），其中 \(\lambda \in K^{*}\) ，就是这种情形。

我们用 \(s\) 表示从 \(\mathrm{E} \oplus \mathrm{F}\) 到 \(\mathrm{F} \oplus \mathrm{E}\) 的希尔伯特空间等距同构，它由 \(s(x, y) = (-y, x)\) （对所有 \((x, y) \in \mathrm{E} \oplus \mathrm{F}\) ）定义。

命题 7. --- 设 u 是从 E 到 F 中的稠定部分算子。

\begin{enumerate}
\def\labelenumi{\alph{enumi})}
\item
  \(u^{*}\) 的图像等于 \(s(\Gamma_{u}^{\circ}) = s(\Gamma_{u})^{\circ}\) ；
\item
  部分算子 \(u^{*}\) 是闭的；
\item
  \(u^{*}\) 的核是 \(u\) 的像的正交补。
\end{enumerate}

记 \(\mathrm{W} = s(\Gamma_u)^\circ\) 。由于线性映射 \(s\) 是酉的，我们有 \(\mathrm{W} = s(\Gamma_u^\circ)\) 。

我们有 \((y,x)\in \mathrm{W}\) 当且仅当

\[
\langle (y, x) \mid (- u (x ^ {\prime}), x ^ {\prime}) \rangle = 0
\]

对所有 \(x' \in \mathrm{dom}(u)\) 成立，亦即若

\[
\langle y \mid u (x ^ {\prime}) \rangle = \langle x \mid x ^ {\prime} \rangle
\]

（对所有 \(x' \in \text{dom}(u)\) ）。当 \(y \in \text{dom}(u^{*})\) 且 \(x = u^{*}(y)\) 时，这一性质成立（参见 p. 236 的公式 (1)）。反之，若这一条件成立，我们推出 \(|\langle y \mid u(x') \rangle| \leqslant \|x\| \|x'\|\) 对所有 \(x' \in \text{dom}(u)\) 成立，这蕴含 \(y\) 属于 \(\text{dom}(u^*)\) ；这时我们有 \(u^*(y) = x\) 。因此 \(W = \Gamma_{u^*}\) 。

算子 \(u^{*}\) 是闭的，因为空间 \(s(\Gamma_{u})^{\circ}\) 在 \(\mathrm{F} \oplus \mathrm{E}\) 中是闭的。

我们来证明断言 \(c\) )。若 \(y\) 与 \(u\) 的像正交，则 D 上的线性形式 \(\lambda_y: x \mapsto \langle y | u(x) \rangle\) 为零，于是 \(y \in \text{dom}(u^*)\) 且 \(u^*(y) = 0\) 。反之，设 \(y \in \text{dom}(u^*)\) 。这时 \(u^*(y) = 0\) 当且仅当 \(y\) 与 \(u(x)\) （对所有 \(x \in D\) ）正交（p. 236 的公式 (1)）。

命题 8. --- 设 u 是从 E 到 F 的稠定算子。则 \(u^{*}\) 稠定当且仅当 u 可闭。在这种情形， u 的闭包 \(\overline{u}\) 等于 \(u^{**}\) ，且 \(\overline{u}\) 的伴随等于 \(u^{*}\) 。

由命题 7，部分算子 \(u^*\) 是闭的。假设 \(u^*\) 的定义域 D* 在 F 中稠密。设 \(u^{**}\) 是 \(u^*\) 的伴随；它是从 E 到 F 的闭部分算子。我们来证明 \(u \subset u^{**}\) ，这将蕴含 \(u\) 可闭。设 \(x \in \text{dom}(u)\) 。按 \(u^*\) 的定义， D* 上由 \(y \mapsto \langle x | u^*(y) \rangle\) 与 \(y \mapsto \langle u(x) | y \rangle\) 所给的线性形式相等；因此我们有 \(x \in \text{dom}(u^{**})\) 且 \(u^{**}(x) = u(x)\) ，由此即得断言。

反之，设 \(u\) 可闭；我们有 \(\Gamma_{\overline{u}} = \overline{\Gamma}_u\) （IV, p. 228 的命题 1）。设 \(y \in F\) 是与 \(\mathrm{dom}(u^*)\) 正交的一个向量。 \((y,0)\) （ \(F \oplus E\) 中的元素）这时属于 \(u^*\) 的图像的正交补。而由命题 7, a) ，我们有

\[
\Gamma_ {u ^ {*}} ^ {\circ} = (s (\Gamma_ {u}) ^ {\circ}) ^ {\circ} = \overline {{s (\Gamma_ {u})}} = s (\overline {{\Gamma}} _ {u})
\]

由此得出 \((0,y)\in\Gamma_{\overline{u}}\) ，从而 \(y=\overline{u}(0)=0\) 。由于 dom \((u^{*})\) 的正交补约化为 0 ，空间 dom \((u^{*})\) 在 F 中稠密。

最后，把命题 7 应用于 \(u^{*}\) ，即得

\[
\Gamma_ {u ^ {* *}} = s ^ {- 1} (\Gamma_ {u ^ {*}} ^ {\circ}) = s ^ {- 1} (s (\Gamma_ {u} ^ {\circ \circ})) = \overline {{\Gamma}} _ {u},
\]

因而 \(u^{**} = \overline{u}\) ，进而 \(\overline{u}^{*} = (u^{*})^{**} = \overline{u^{*}} = u^{*}\) ，因为 \(u^{*}\) 是闭的。

推论. --- 若 u 是从 E 到 F 中的闭稠定部分算子，则 \(u^{*}\) 稠定，且有 \(u^{**} = u\) 。

定义 7. --- 设 u 是 E 上的部分算子。若 u 稠定且 \(u^{*}\) 是 u 的扩张，就称 u 是对称的。若 u 稠定且 \(u^{*} = u\) ，就称 u 是自伴的。

称 u 是本质自伴的，如果它可闭且其闭包 \(\overline{u}\) 是自伴的。

称 u 是下有界部分算子，如果 u 是对称的，且存在实数 c 使得 \(\langle x | u(x) \rangle \geqslant c \|x\|^{2}\) 对属于 u 的定义域的所有 x 成立。这时称 c 是 u 的一个下界。若 c = 0 ，也称 u 是正部分算子。

我们用 \(\mathcal{A}(\mathrm{E})\) 表示 E 上自伴部分算子所成的集。

注. --- 1) 算子 \(u\) 在 E 上稠定且对称，当且仅当下式成立：

\[
\langle x \mid u (y) \rangle = \langle u (x) \mid y \rangle\tag{2}
\]

（对所有 \((x,y)\in\mathrm{dom}(u)^{2}\) ）。这一公式表明 u 的定义域包含在 \(u^{*}\) 的定义域中，进而 \(u^{*}\) 与 u 在 u 的定义域上重合。特别地，由此得出 \(\langle x|u(x)\rangle\in\mathbf{R}\) 对所有 \(x\in\mathrm{dom}(u)\) 成立。

正如在多个例子中将看到的那样，公式 (2) 常常可以通过直接计算加以验证。另一方面，伴随的定义域的精确确定——唯有它才能判别一个对称算子是否自伴——可能是非常精细的。

\begin{enumerate}
\def\labelenumi{\arabic{enumi})}
\setcounter{enumi}{1}
\item
  自伴部分算子 u 是本质自伴的（参见命题 8）。
\item
  设 \(u\) 是 E 上的对称部分算子。算子 \(u\) 可闭（命题 7, b)）。它满足 \(\text{dom}(u) \subset \text{dom}(u^*)\) ，且 \(u\) 自伴当且仅当 \(\text{dom}(u) = \text{dom}(u^*)\) 。此外， \(\overline{u}\) 的闭包 \(u\) 是对称的，因为 \(\overline{u} \subset u^* = \overline{u}^*\) （命题 8）。
\item
  设 K = C 。设 \(u \in \mathcal{P}(\mathrm{E};\mathrm{E})\) 是稠定部分算子。条件 \(\langle x | u(x) \rangle \in \mathbf{R}\) （对所有 \(x \in \operatorname{dom}(u)\) ）蕴含 u 是对称的（EVT, V, p. 2，注）；特别地，若 \(\langle x | u(x) \rangle \in \mathbf{R}_{+}\) 对所有 \(x \in \operatorname{dom}(u)\) 成立，则 u 是正的。
\item
  设 u 与 v 是 E 上的对称部分算子。若 u 自伴且 \(u \subset v\) ，则 \(v \subset v^{*} \subset u^{*} = u\) ，因此 u = v 。
\item
  本质自伴部分算子 \(u\) 是对称的，因为 \(u \subset \overline{u}\) 蕴含 \(\overline{u} = \overline{u}^{*} \subset u^{*}\) ，从而 \(u \subset u^{*}\) 。
\item
  设 u 与 v 是 E 上的对称部分算子。若 \(u+v\) 稠定（例如 u 或 v 属于 \(\mathcal{L}(\mathrm{E})\) ），则 \(u+v\) 是对称的。一般而言，部分算子 \(u + v\) 不是自伴的，即使 u 与 v 都自伴（IV, p. 347 的习题 9）。
\item
  设 \(u\) 是 E 上的对称部分算子。实数 \(c\) 是 \(u\) 的下界，当且仅当算子 \(u - c \cdot 1_{\mathrm{E}}\) 是正的。
\end{enumerate}

引理 3. --- 假设 K = R 。设 u 是从 E 到 F 中的稠定部分算子。

\begin{enumerate}
\def\labelenumi{\alph{enumi})}
\item
\(u_{(\mathbf C)}\) 的伴随是 \((u^*)_{(\mathbf C)}\) ；
\item
  设 E = F ；部分算子 u 是对称的（相应地，自伴的），当且仅当部分算子 \(u_{(\mathbf{C})}\) 是对称的（相应地，自伴的）。
\end{enumerate}

我们来证明 \(a\) )。设 \(y \in \mathrm{F}_{(\mathbf{C})}\) ，并把它写成 \(y = y_1 + iy_2\) ，其中 \(y_1, y_2 \in \mathrm{F}\) 。对所有 \((x_1, x_2) \in \mathrm{E} \times \mathrm{E}\) ，我们有

\[
\begin{array}{c} \langle u _ {(\mathbf {C})} (x _ {1} + i x _ {2})   |   y \rangle = \langle u (x _ {1})   |   y _ {1} \rangle + i \langle u (x _ {1})   |   y _ {2} \rangle \\ - i \langle u (x _ {2})   |   y _ {1} \rangle + \langle u (x _ {2})   |   y _ {2} \rangle . \end{array}
\]

若 \(y \in \text{dom}(u^*)(\mathbf{C})\) ，我们推出 \(y \in \text{dom}((u_{(\mathbf{C})})^*)\) 且 \(u_{(\mathbf{C})}^*(y) = (u_{(\mathbf{C})})^*(y)\) ，从而 \(u_{(\mathbf{C})}^* \subset (u_{(\mathbf{C})})^*\) 。

反之，设 \(y\) 属于 \(\text{dom}((u_{(\mathbf{C})})^*)\) 。在上式中取 \(x_2 = 0\) （相应地， \(x_1 = 0\) ），即可验证 \(y_1 \in \text{dom}(u^*)\) （相应地， \(y_2 \in \text{dom}(u^*)\) ），从而 \(y \in \text{dom}(u^*)(\mathbf{C})\) 。

断言 \(a\) ) 蕴含：若 \(u_{(\mathbf{C})}\) 是对称的（相应地，自伴的），则 \(u\) 也是对称的（相应地，自伴的）。

反之，设 \(u_{(\mathbf{C})}\) 是对称的。关系式 \(\langle u(x)|y\rangle = \langle x|u(y)\rangle\) （对所有 \((x,y) \in \mathrm{dom}(u_{(\mathbf{C})}) \times \mathrm{dom}(u_{(\mathbf{C})})\) ）蕴含：在取 x 与 y 属于 \(\mathrm{dom}(u)\) （ \(\mathrm{dom}(u_{(\mathbf{C})})\) 的子空间）时， u 是对称的。若 \(u_{(\mathbf{C})}\) 自伴，则前述论证表明 u 是对称的；由于 \(\mathrm{dom}(u^{*}) = \mathrm{dom}(u_{(\mathbf{C})}^{*}) \cap \mathrm{F}\) ，断言 a) 蕴含 \(\mathrm{dom}(u^{*}) = \mathrm{dom}(u_{(\mathbf{C})}) \cap \mathrm{F} = \mathrm{dom}(u)\) ，因此 u 是自伴的。

\subsection{自伴算子的初等判别法}

命题 9. --- 设 \(v \in \mathcal{L}(\mathrm{F};\mathrm{E})\) 是从 F 到 E 中的连续单射线性映射，其像在 E 中稠密。 \(v\) 的伴随是从 E 到 F 中的连续单射线性映射，且有 \((v^{*})^{-1} = (v^{-1})^{*}\) 。特别地，若 E = F ，则自同态 \(v\) 是埃尔米特的，当且仅当部分算子 \(v^{-1}\) 是自伴的。

部分算子 \(v^{-1}\) 是从 E 到 F 中的闭稠定算子（IV, p. 228 的例 2），而 v 的伴随 \(v^{*}\) 是从 E 到 F 中的连续线性映射；它是单射的，因为 v 的像在 E 中稠密（EVT, V, p. 41，命题 4）。设 s （相应地， \(s'\) ）是从 E ⊕ F 到 F ⊕ E 上的等距同构 \((x,y)\mapsto(-y,x)\) （相应地，从 F ⊕ E 到 E ⊕ F 上的等距同构 \((y,x)\mapsto(-x,y)\) ），并设 \(\iota\) （相应地， \(\iota'\) ）是从 F ⊕ E 到 E ⊕ F 上的等距同构 \((y,x)\mapsto(x,y)\) （相应地，从 E ⊕ F 到 F ⊕ E 上的等距同构 \((x,y)\mapsto(y,x)\) ）。这时有 \(s\circ\iota=-\iota'\circ s'\) ，由此

\[
\Gamma_ {(v ^ {- 1}) ^ {*}} = s \left(\Gamma_ {v ^ {- 1}}\right) ^ {\circ} = s \left(\iota \left(\Gamma_ {v}\right)\right) ^ {\circ} = - \iota^ {\prime} \left(s ^ {\prime} \left(\Gamma_ {v}\right)\right) ^ {\circ} = - \iota^ {\prime} \left(\Gamma_ {v ^ {*}}\right) = \Gamma_ {(v ^ {*}) ^ {- 1}}
\]

（依据 IV, p. 236 的命题 7）。命题得证。

命题 10（海林格–特普利茨）. --- 设 u 是希尔伯特空间 E 上的对称部分算子。若 u 的定义域等于 E ，则 \(u \in \mathcal{L}(E)\) ，且 u 是埃尔米特的。

事实上，部分算子 u 是可闭的（IV, p. 237 的命题 8），而由于它的定义域是 E ，它是闭的（IV, p. 228，注）。于是引用 EVT, I, p. 19，推论 5 即得结论。

推论. --- 设 u 是 E 上的对称部分算子。若部分算子 u 诱导一个从 dom(u) 到 E 上的双射线性映射，则 u 是自伴的。

事实上， u 的逆部分算子 \(u^{-1}\) 是对称的且以 E 为定义域（命题 9），于是 \(u^{-1}\) 是 \(\mathcal{L}(\mathrm{E})\) 中的自伴元素（命题 10），从而 u 是埃尔米特的（命题 9）。

命题 11. --- 设 u 是 E 上的对称部分算子，并设 \(\lambda \in \mathbf{C}\) 。若 \(u + \lambda 1_{\mathrm{E}}\) 与 \(u + \overline{\lambda} 1_{\mathrm{E}}\) 都是满射的，则 u 是自伴的。

只须证明 \(\operatorname{dom}(u^{*}) \subset \operatorname{dom}(u)\) 。设 \(x \in \operatorname{dom}(u^{*})\) 。由假设，存在 \(y \in \operatorname{dom}(u)\) 使得 \(u(y) + \overline{\lambda}y = u^{*}(x) + \overline{\lambda}x\) 。我们来证明 y = x 。对所有 \(z \in \operatorname{dom}(u)\) ，得到

\[
\begin{array}{r l} & {\langle (u + \lambda 1 _ {\mathrm{E}}) (z) | x \rangle = \langle z | (u ^ {*} + \overline {{\lambda}} 1 _ {\mathrm{E}}) (x) \rangle} \\ & {\qquad = \langle z | (u + \overline {{\lambda}} 1 _ {\mathrm{E}}) (y) \rangle = \langle (u + \lambda 1 _ {\mathrm{E}}) (z) | y \rangle ,} \end{array}
\]

（因为 u 是对称的）。由于算子 \(u + \lambda1_{E}\) 是满射的，确有 y = x ，从而 \(x \in \text{dom}(u)\) 。

我们稍后将看到（参见 IV, p. 248 的命题 17）：若 \(\lambda \in \mathbf{C} - \mathbf{R}\) ，则逆命题成立。

命题 12. --- 设 u 是从 E 到 F 中的闭稠定算子。 E 上的部分算子 \(u^{*}\circ u\) 是自伴且正的。它的定义域是 u 的一个核。

用 \(v\) 表示部分算子 \(1_{\mathrm{E}} + u^{*} \circ u\) 。它的定义域是 \(\operatorname{dom}(u^{*} \circ u)\) ，包含在 \(\operatorname{dom}(u)\) 中。对所有 \(x \in \operatorname{dom}(u)\) 与 \(y \in \operatorname{dom}(v)\) ，我们有 \(u(y) \in \operatorname{dom}(u^{*})\) 以及

\begin{enumerate}
\def\labelenumi{(\arabic{enumi})}
\setcounter{enumi}{2}
\tightlist
\item
\end{enumerate}

\[
\langle x \mid v (y) \rangle = \langle x \mid y \rangle + \langle x \mid (u ^ {*} \circ u) (y) \rangle = \langle x \mid y \rangle + \langle u (x) \mid u (y) \rangle ,\tag{4}
\]

\[
\langle y \mid v (y) \rangle = \| y \| ^ {2} + \| u (y) \| ^ {2}.
\]

公式 (4) 蕴含部分算子 \(v\) 是单射的。此外，由 IV, p. 236 的命题 7 \(a\) )，有 \(\mathrm{F} \oplus \mathrm{E} = \Gamma_{u^*} \oplus s(\Gamma_u)\) 。设 \(x \in \mathrm{E}\) 。存在 \(y' \in \mathrm{dom}(u^*)\) 与 \(x' \in \mathrm{dom}(u)\) 使得

\[
(0, x) = (y ^ {\prime}, u ^ {*} (y ^ {\prime})) + (- u (x ^ {\prime}), x ^ {\prime}) = (y ^ {\prime} - u (x ^ {\prime}), x ^ {\prime} + u ^ {*} (y ^ {\prime})).
\]

因此我们有 \(y' = u(x')\) ，从而 \(x' \in \text{dom}(u^* \circ u) = \text{dom}(v)\) 以及

\[
x = x ^ {\prime} + u ^ {*} (y ^ {\prime}) = x ^ {\prime} + (u ^ {*} \circ u) (x ^ {\prime}) = v (x ^ {\prime}).
\]

因此， E 上的部分算子 v 是满射的，并诱导一个从 dom(v) 到 E 上的双射线性映射。

设 \(x \in \mathrm{E}\) 与 \(\operatorname{dom}(v)\) 正交。把 x 写成 \(x = v(x')\) ，其中 \(x' \in \operatorname{dom}(v)\) 。由公式 (4) ，我们得到

\[
0 = \langle x ^ {\prime} | x \rangle = \langle x ^ {\prime} | v (x ^ {\prime}) \rangle = \| x ^ {\prime} \| ^ {2} + \| u (x ^ {\prime}) \| ^ {2}
\]

由此 \(x' = 0\) ，进而 \(x = 0\) 。因此 \(v\) 的定义域在 E 中稠密。

部分算子 \(v\) 稠定；它是双射的，且公式 (3) 表明它是对称的。对 \(v\) 应用命题 10 的推论，我们断定它是自伴的。于是 \(u^{*} \circ u = v - 1_{\mathrm{E}}\) 是自伴的。此外，公式

\[
\langle x \mid (u ^ {*} \circ u) (x) \rangle = \| u (x) \| ^ {2}
\]

（对所有 \(x \in \text{dom}(u^{*} \circ u)\) ）蕴含 \(u^{*} \circ u\) 是正的。

最后，设 \(y \in \mathrm{E}_u\) 与 \(\operatorname{dom}(u^* \circ u)\) 正交。存在 \(x\) 的定义域中的元素 \(u^* \circ u\) 使得 \(y = v(x) = x + (u^* \circ u)(x)\) 。这时我们有 \(0 = (x \mid y)_u = \langle y \mid y \rangle\) ，由此 \(y = 0\) 。

\subsection{微分算子}

设 \(n \in N\) ， U 是 \(R^{n}\) 的开子集。我们给 \(R^{n}\) 与 U 赋以勒贝格测度，记作 \(\mu\) 。

设 \(k \in N\) 与 \(h \in N\) 满足 \(h \geqslant k\) 。设 D 是 U 上阶 \(\leqslant k\) 的纯量微分算子，其系数 \((n_{\alpha})_{|\alpha| \leqslant k}\) 在 U 上属于类 \(C^{h}\) 。假设对每个满足 \(\alpha\) 的 \(|\alpha| \leqslant k\) 以及每个满足 \(\beta\) 的 \(0 \leqslant \beta \leqslant \alpha\) ，函数 \(\partial^{\beta} n_{\alpha}\) 在 U 上有界。

设 \({}^{t}D\) 是 D 的转置纯量微分算子（VAR, R2, 14.3.2）；它的阶 \(\leqslant k\) ，属于类 \(C^{h-k}\) ；对 \(\varphi\in\mathcal{D}(U)\) ，我们有

\[
{ } ^ { t } \mathrm{D} ( \varphi ) = \sum _ { | \alpha | \leqslant k } ( - 1 ) ^ { | \alpha | } \partial ^ { \alpha } ( \overline { { n } } _ { \alpha } \varphi )
\]

（loc. cit.）；特别地， \({}^{t}D\) 的系数在 U 上有界。

我们用 D\_ 表示 \(L^{2}(U)\) 上以 \(\mathcal{D}(U)\) 为定义域、由下式定义的部分算子：

\[
\varphi \mapsto \mathrm{D} (\varphi) = \sum_ {| \alpha | \leqslant k} n _ {\alpha} \partial^ {\alpha} \varphi .\tag{5}
\]

设 \(H_{D}\) 是这样的 \(f \in L^{2}(U)\) 所成的空间：下列广义函数

\[
\mathrm{D} (f) = \sum_ {| \alpha | \leqslant k} n _ {\alpha} \partial^ {\alpha} f
\]

属于 \(L^{2}(U)\) ；我们用 \(D_{+}\) 表示以 \(H_{D}\) 为定义域、由 \(f \mapsto \mathrm{D}(f)\) 所定义的部分算子。

由于我们有 \(\partial^{\alpha}f\in \mathrm{L}^{2}(\mathrm{U})\) （若 \(f\in \mathrm{H}^k (\mathrm{U})\) 且 \(|\alpha |\leqslant k\) ），索伯列夫空间 \(\mathrm{H}^k (\mathrm{U})\) 包含在 \(\mathrm{H}_{\mathrm{D}}\) 中；一般而言，这两个空间是不同的。

我们有 \(D_{-} \subset D_{+}\) ，它们都是 D 在 \(L^{2}(U)\) 上的相伴微分算子（IV, p. 235 的定义 5）。

命题 13. --- 设 u 是 L²(U) 上的部分算子。若 D₋ ⊂ u ⊂ D₊ ，则 u 可闭，且 (ᵗD)₋ ⊂ u* ⊂ (ᵗD)₊ 。

设 \(\varphi\) 与 \(\psi\) 属于 \(\mathcal{D}(\mathrm{U})\) 。这时我们有

\[
\begin{array}{c} \langle \varphi | \mathrm{D} (\psi) \rangle = \sum_ {| \alpha | \leqslant k} \int_ {\mathrm{U}} n _ {\alpha} \overline {{\varphi}} \partial^ {\alpha} \psi d \mu \\ = \sum_ {| \alpha | \leqslant k} (- 1) ^ {| \alpha |} \langle \partial^ {\alpha} (\overline {{n}} _ {\alpha} \varphi) | \psi \rangle = \langle {} ^ {t} \mathrm{D} (\varphi) | \psi \rangle \end{array}
\]

（参见 VAR, R2, 14.3.8）。由于 \(\mathcal{D}(\mathrm{U})\) 在 \(\mathrm{L}^2 (\mathrm{U})\) 中稠密，这蕴含 \(\varphi \in \mathrm{dom}(u^{*})\) 且 \(u^{*}(\varphi) = {}^{t}\mathrm{D}(\varphi)\) 。因此我们有 \(({}^{t}\mathrm{D})_{-}\subset u^{*}\) 。特别地， \(u^{*}\) 稠定而 \(u\) 可闭（IV, p. 237 的命题 8）。

设 \(f \in \mathrm{dom}(u^{*})\) 与 \(\varphi \in \mathcal{D}(\mathrm{U})\) 。由于 \(\mathcal{D}(\mathrm{U}) \subset \mathrm{dom}(u)\) ，与 \(u^{*}(f)\) 相伴的广义函数满足

\[
\langle u ^ {*} (f), \varphi \rangle = \langle \overline {{\varphi}} | u ^ {*} (f) \rangle = \langle u (\overline {{\varphi}}) | f \rangle .
\]

由于 \(D_{-} \subset u\) ，可由公式 (5) 计算 \(u(\overline{\varphi})\) ，由此

\[
\begin{array}{l} \langle u ^ {*} (f), \varphi \rangle = \sum_ {\alpha} \langle n _ {\alpha} \partial^ {\alpha} \overline {{\varphi}} | f \rangle = \sum_ {\alpha} \langle \partial^ {\alpha} \overline {{\varphi}} | \overline {{n}} _ {\alpha} f \rangle \\ \qquad = \sum_ {\alpha} \langle \overline {{n}} _ {\alpha} f, \partial^ {\alpha} \varphi \rangle = \sum_ {\alpha} (- 1) ^ {| \alpha |} \langle \partial^ {\alpha} (\overline {{n}} _ {\alpha} f), \varphi \rangle = \langle {} ^ {t} D (f), \varphi \rangle . \end{array}
\]

于是广义函数 \(u^{*}(f)\) 与 \({}^{t}\mathrm{D}(f)\) 相等；因而广义函数 \(f\) 属于 \(H_{tD}\) ，且 \(u^{*}(f)={}^{t}\mathrm{D}(f)\) ，由此 \(u^{*}\subset({}^{t}\mathrm{D})_{+}\) 。

注. --- 这一命题表明，部分算子 \(u\) （它满足 \(\mathrm{D}_{-} \subset u \subset \mathrm{D}_{+}\) ）的伴随总可以在广义函数的意义下计算：元素 \(f\) （它们属于 \(u^{*}\) 的定义域）是 \(\mathrm{L}^{2}(\mathrm{U})\) 中这样的元素：广义函数 \(^{t}\mathrm{D}(f)\) 属于 \(\mathrm{L}^{2}(\mathrm{U})\) ，且我们有 \(u^{*}(f) = ^{t}\mathrm{D}(f)\) 。

若作为纯量微分算子有 \({}^{t}D = D\) ，就称 D 是形式对称的。在这种情形，部分算子 D\_ 是对称的。

考虑由下式定义的二阶纯量微分算子的特殊情形：

\[
\Delta = - \sum_ {i = 1} ^ {n} \partial_ {i} ^ {2}.
\]

拉普拉斯算子（ U 上的），指的是 \(L^{2}(U)\) 上自伴的部分算子 u ，满足 \(\Delta_{-} \subset u\) （参见 VAR, R2, 14.4.3, p. 83）。我们稍后将看到（IV, p. 261，例）， \(L^{2}(U)\) 上总存在至少一个拉普拉斯算子；也可能存在不止一个（IV, p. 358 的习题 17）。